% -----------------------------------------------------------
\section{1830}
% -----------------------------------------------------------
	\label{EMS-1830}
	
	% ----------------------
	\subsection{Joseph Smith}
	% ----------------------
		\label{EMS-1830-JS}

		% -----
		\subsubsection{Introduction to 1830}
		% -----
			\label{EMS-1830-JS-INTRO}
		
			Include more Moses?
			
		% -----
		\subsubsection{Revelation, circa Early 1830}
		% -----
			\label{EMS-1830-JS-CIR-EARLY-1830}
		
			% --
			\paragraph{Sources and Access}
			% --

				\begin{tabular}{p{1in}p{1.5in}p{1in}p{1.5in}}
				JSP	 	& D1:108--112	& JSR & 372--374 \\
				DC	 	& ---		 	& HC 		& --- \\
				HDDC 	& ---			& TCHC 		& --- \\
				JSUR	& 32--35		& Online  & \href{https://www.josephsmithpapers.org/paper-summary/revelation-circa-early-1830/1}{JSP Website} \\
				\end{tabular}
			
			% --
			\paragraph{Original Source and Background}
			% --
			
				The source is RB1: ``Revelation, Manchester Township, Ontario Co., NY, [ca. early] 1830. Featured version, titled `23 Commandment AD 1830,' copied [ca. Mar. 1831] in Revelation Book 1, pp. 30--31; handwriting of John Whitmer; CHL. Includes redactions.''
				
				Note that in the JSR part of the table above, Marquardt does not give the text of this revelation, not knowing at the time whether it still existed. His discussion still provides useful historical background.
				
				From the JSP: ``This revelation directed JS and others to obtain the copyright for the Book of Mormon `upon all the face of the Earth' and appointed four men, Oliver Cowdery, Joseph Knight Sr., Hiram Page, and Josiah Stowell, to travel to Kingston, Upper Canada, to secure a copyright for the four provinces of Canada. Since the men apparently traveled to Kingston sometime between late January and early March 1830, the revelation must have been written early in 1830. More than six months earlier, on 11 June 1829, JS had applied for a copyright for the Book of Mormon in the United States by submitting the book's title page to the clerk of the U.S. District Court for the Northern District of New York. The copyright protected the text against unauthorized changes and distribution. In January 1830, Abner Cole, who printed the Palmyra \textit{Reflector} on the same press being used for the Book of Mormon, tested the prepublication copyright by printing a portion of the Book of Mormon in his newspaper without authorization. The problem was quickly resolved, in part because JS asserted his copyright authority. Shortly after that confrontation, this revelation commanded JS to expand copyright protection of the Book of Mormon throughout the world.
				
				``The potential profits from the sale of the books had also recently become the subject of an agreement between JS and Martin Harris, the primary financier for the printing, who had mortgaged his property in August 1829 as payment to E. B. Grandin. On 16 January 1830, likely during his visit to Palmyra to confront Cole, JS signed an agreement that granted Harris the right to sell copies of the Book of Mormon until he was compensated for the cost of the printing. Although the agreement permitted Harris to collect \$3,000 from the sale of the books, this revelation, possibly dictated about the same time, explicitly excluded Harris from sharing in the temporal blessings associated with the wider sale and distribution of the Book of Mormon. As shown in the text transcribed below, the revelation promised blessings to `those who have assisted' JS, `yea even all save M$\diamond\diamond$tin [Martin] only.' Three layers of deletion obscured Harris's name. Hiram Page later wrote that the necessary preparations for the trip to Upper Canada were made `in a sly manor so as to keep Martin Haris from drawing a share of the money.'
				
				``According to a much later recollection by David Whitmer, Hyrum Smith originally suggested to JS that the Book of Mormon copyright could be sold for `considerable money' in Upper Canada. Hiram Page recalled that a small group of church leaders were assembled at the Smith log home in Palmyra Township when this revelation came. It directed Cowdery, Page, Stowell, and Knight to obtain the copyright for the Book of Mormon not only in Canada but `upon all the face of the Earth of which is known by you.' They were to begin their efforts in Kingston and were then authorized, under certain conditions, to sell the Book of Mormon copyright for the four major British provinces in North America.
				
				``Cowdery, Page, and likely the other two traveled to Kingston in response to the revelation, but they failed to obtain a copyright. A local resident apparently informed them that they were more likely to secure a copyright if they went to York, the capital of Upper Canada, but rather than travel more than 160 miles farther, they returned home and reported to JS at the home of Peter Whitmer Sr. in Fayette. Almost six decades later David Whitmer questioned the legitimacy of the revelation because of the group's failure to secure a copyright. He claimed that in a council held on 1 November 1831 to consider printing the revelations, JS repudiated the revelation.%
					\footnote{%
						% \label{}%
						JSP footnote: ``Traughber, `False Prophecies,' 1--2; Whitmer, \textit{Address to All Believers in Christ}, 31. The extant minutes of the 1 November 1831 conference do not contain any discussion of this revelation.''
						}
				Whitmer further stated that he and Jacob and John Whitmer were present when Cowdery and Page returned. `We asked Joseph,' he wrote, `how it was that he had received a revelation from the Lord for some brethren to go to Toronto [Kingston] and sell the copy-right, and the brethren had utterly failed in their undertaking.' Joseph reportedly enquired of the Lord through the Urim and Thummim and, according to Whitmer, received a revelation that stated, `Some revelations are of God: some revelations are of man: and some revelations are of the devil.'%
					\footnote{%
						% \label{}%
						JSP footnote: ``Whitmer, \textit{Address to All Believers in Christ}, 31; see also Traughber, `False Prophecies,' 2.''
						}
				Page, however, left no indication in his recollection that he was bitter about the revelation or his trip to Canada.%
					\footnote{%
						% \label{}%
						JSP footnote: ``H. Page to W. McLellin, 2 Feb. 1848. Page's letter to McLellin offered this comment on his and Cowdery’s trip to Kingston: `But when we got there, there was no purchaser neither were they authorized at Kingston to buy rights for the provence; but little York was the place where such buisaness had to be done, we were to get 8,000 dollars[.] we were treated with the best of respects by all we met with in Kingston---by the above we may learn how a revlation may be received and the person receving it not be benafited.'''
						}
				Even though the text in Revelation Book 1 includes editing marks, it was not published with other revelations in the Book of Commandments in 1833.''
			
			% --
			\paragraph{Text}
			% --
			
				\begin{tightcenter}
					\hypertarget{EMS-1830-JS-CIR-EARLY-1830-a}%
				23
					\marginnote{a}
				Commandment AD 1830
				\end{tightcenter}
				
				A Revelation given to Joseph Oliver [Cowdery] Hyram [Hiram Page] Josiah [Stowell] \& Joseph Knight [Sr.] given at Manchester Ontario C[ounty] New York
				
				Behold I the Lord am God I Created the Heavens \& the Earth \& all things that in them is wherefore they are mine \& I sway my scepter over all the Earth \& ye are in my hands to will \& to do that I can deliver you out of evry difficulty \& affliction according to your faith \& dilligence \& uprightness Before me
					\hypertarget{EMS-1830-JS-CIR-EARLY-1830-b}%
				\phantom{ }
					\marginnote{b}%
				\phantom{ }\& I have covenanted with my Servent that earth nor Hell combined against him shall not take the Blessing out of his hands which I have prepared for him if he walketh uprightly before me neither the spiritual nor the temporal Blessing \& Behold I also covenanted with those who have assisted him in my work that I will do unto them even the same Because they have done that which is pleasing in my sight (yea even all save \sout{M$\diamond\diamond$tin} \sout{only} it be one only)
					\hypertarget{EMS-1830-JS-CIR-EARLY-1830-b}%
				Wherefore be dilligent in Securing the Copy
					\marginnote{b}%
				right of my \sout{Servent} work upon all the face of the Earth of which is known by you unto unto my Servent Joseph \& unto him whom he willeth accordinng as I shall command him that the faithful \& the righteous may retain the temperal Blessing as well as the Spirit[u]al \& also that my work be not destroyed by the workers of iniquity to their own distruction \& damnation when they are fully ripe \& now Behold I say unto you that I have covenanted
					\hypertarget{EMS-1830-JS-CIR-EARLY-1830-c}%
				\phantom{ }
					\marginnote{c}%
				\phantom{ }\& it Pleaseth me that Oliver Cowderey Joseph Knight Hyram Page \& Josiah Stowel shall do my work in this thing yea even in securing the <Copy> right \& they shall do it with an eye single to my Glory that it may be the means of bringing souls unto \sout{me} Salvation through mine only Begotten Behold I am God I have spoken it \& it is expedient in me Wherefor I say unto you that ye shall go to Kingston seeking me continually through mine only Begotten \& if ye do this ye shall have my spirit to go with you \& ye shall have an addition of all things which is expedient in me.
					\hypertarget{EMS-1830-JS-CIR-EARLY-1830-d}%
				\phantom{ }
					\marginnote{d}%
				\phantom{ }\& I grant unto my servent a privelige that he may sell <a copyright> through you speaking after the manner of men for the four Provinces if the People harden not their hearts against the enticeings of my spirit \& my word for Behold it lieth in themselves to their condemnation \sout{\&} or to their salvation
					\hypertarget{EMS-1830-JS-CIR-EARLY-1830-e}%
				Behold
					\marginnote{e}%
				my way is before you \& the means I will prepare \& the Blessing I hold in mine own hand \& if ye are faithful I will pour out upon you even as much as ye are able to Bear \& thus it shall be Behold I am the father \& it is through mine only begotten which is Jesus Christ your Redeemer amen
				
		% -----
		\subsubsection{Articles and Covenants, circa April 1830 [D\&C 20]}
		% -----
			\label{EMS-1830-JS-CIR-APR-1830}
		
			% --
			\paragraph{Sources and Access}
			% --

				\begin{tabular}{p{1in}p{1.5in}p{1in}p{1.5in}}
				JSP	 	& D1:116--126	& JSR 		& 62--68	 \\
				DC	 	& Section 20	& HC 		& 1:64--70	 \\
				HDDC 	& 286--351		& TCHC 		& 1:47--51	 \\
				Online  & \href{https://www.josephsmithpapers.org/paper-summary/articles-and-covenants-circa-april-1830-dc-20/1}{JSP Website} & 
				JWGB	& \hyperref[DC20]{DC Section 20} \\
				\end{tabular}
			
			% --
			\paragraph{Original Source and Background}
			% --
			
				The source is a newspaper article from 1831: ```Articles and Covenants,' [Fayette Township, Seneca Co., NY, ca. Apr. 1830, though parts may have been received as early as ca. summer 1829]. Featured version part of `The Mormon Creed,' in \textit{Painesville (OH) Telegraph}, 19 Apr. 1831, vol. 2, no. 44 (second series), p. [4]. The microfilm copy of the text transcribed herein was filmed by the Microfilm Corporation of Cleveland, OH, 1947, copy at CHL. The \textit{Painesville Telegraph} version and the copy found in Revelation Book 1 both appear to have been created about the same time, but differences between the two versions indicate that the former was based on an earlier copy; therefore, the \textit{Telegraph} version is featured here.''
				
				For the background, see \hyperref[DC20-HB]{here}.
			
			% --
			\paragraph{Text}
			% --
			
				The articles and covenants of the Church of Christ agreeable to the will and commandments of God. The rise of the Church of Christ in these last days, being 1830 years since the coming of our Lord and Saviour Jesus Christ in the flesh, it being regularly organized and established agreeable to the laws of our country, by the will and commandments of God, in the 4th month, and on the 6th day of the same, which commandments were given to Joseph Smith, jun. who was called of God and ordained an apostle of Jesus Christ, an elder of the church, and also to Oliver Cowdery, who was also called of God an apostle of Jesus Christ, an elder of the church, and ordained under his hand, and this according to the grace of God the Father, and our Lord Jesus Christ, to whom be all glory both now and ever---amen.
				
				For after that it truly was manifested unto the first elder that he had received remission of his sins, he was entangled again in the vanities of the world, but after truly repenting, God visited him by an holy angel, whose countenance was as lightning, and whose garments were pure and white above all whiteness, and gave unto him commandments which inspired him from on high, and gave unto him power, by the means of which was before prepared that he should translate a book; which book contains a record of a fallen people, and also the fullness of the gospel of Jesus Christ to the Gentiles and also to the Jews, proving unto them that the holy scriptures be true, and also that God doth inspire men and call them to his holy work in these last days, as well as in days of old, that he may be the same God forever---amen.

				Which book, given by inspiration, is called the Book of Mormon, and is confirmed to others by the ministering of angels, and declared unto the world by them. Wherefore, having so great witnesses, by them shall the world be judged, even as many as shall hereafter receive this work, either to faith and righteousness, or to the hardness of heart in unbelief to their own condemnation. For the Lord God hath spoken it, for we elders of the church have heard and bear record to the words of the glorious Majesty on high, to whom be glory forever and ever---Amen. Wherefore, by these things, we know that there is a God in heaven, who is infinite and eternal, from everlasting to everlasting the same unchangeable God. the Maker of heaven and earth, and all things that in them is, and that he is all power, and all wisdom, and all understanding, and that he created man male and female after his own image and in his own likeness created he them, and that he gave unto the children of men a commandment that they should love and serve him the only being whom they should worship; but by the transgression of these holy laws, men became sensual and devlish, and became fallen man---wherefore the Almighty God gave his only begotten Son, as is written in those scriptures which hath been given of him, that he suffered temptations but gave no heed unto them, that he was crucified and died and rose again the third day, and that he ascended into heaven to sit down on the right hand of the Father, to reign with almighty power according to the will of the Father, that as many as would believe and were baptized into his holy name and endured in faith to the end should be saved; yea, even as many as were before he came in the flesh, from the beginning, which believed in the words of the holy prophets, which were inspired by the gift of the Holy Ghost, which truly testified of him in all things, as well they which should come after, which should believe in the gifts and calling of God by the Holy Ghost which beareth record of the Father and of the Son, which Father and the Son, and the Holy Ghost is one God, infinite and eternal, without end.--- Amen.

				And we know that all men must repent and believe on the name of Jesus Christ, and worship the Father in his name, and endure in faith on his name to the end, or they cannot be saved in the kingdom of God. And we know that justification through the grace of our Lord and Saviour Jesus Christ is just and true. And we also know that sanctification through the grace of our Lord and Saviour Jesus Christ is just and true to all those who love and serve God with all their mights, minds, and strength. But there is a possibility that men may fall from grace, and depart from the living God, therefore let the church take heed, and pray always, lest they enter into temptations; yea, and he that is sanctified also. And we know that these things are true and agreeable to the revelations of Jesus Christ which was signified by his angel unto John, neither adding nor diminishing to the prophecy of his book; neither to the holy scriptures; neither to the revelations of God which shall come hereafter by the gift and power of the Holy Ghost, neither by the voice of God, neither by the ministering of angels. And the Lord God hath spoken it---and honor, power, and glory be rendered to his holy name both now and ever.---Amen.

				And again, by way of commandment to the church concerning the manner of baptism, behold whosoever humbleth himself before God and desireth to be baptized, and comes forth with a broken heart and a contrite spirit, and witnesseth unto the church that they truly repent of all their sins, and are willing to take upon them the name of Christ, having a determination to serve him unto the end, and truly manifest by their works that they have received the gift of Christ unto the remission of their sins, then shall they be received unto baptism into the church of Christ.

				The duty of Elders, Priests, Teachers, and Deacons, and members of the church of Christ:---An apostle is an elder, and it is his calling to baptize and to ordain other elders, priests, teachers, and deacons, and to administer the flesh and blood of Christ according to the scriptures, and to teach, expound, and exhort, and to baptize and to watch over the church, and to confirm the church by the laying on of hands and the giving of the Holy Ghost, and to take the lead of all meetings, \&c.

				The elders are to conduct the meetings as they are led by the Holy Ghost.

				The priests' duty is to preach, teach, expound, and exhort, and baptize, and administer the sacrament, and visit the house of each member, and exhort them to pray vocally and in secret, and also to attend to all family duties, to ordain priests, teachers, and deacons, and to take the lead in meetings; but none of these offices is he to do when there is an elder present, but in all cases is to assist the elder, \&c. The teacher's duty is to watch over the church always, and be with them and strengthen them, and see that there is no iniquity in the church, nor no hardness with each other, nor no lying nor backbiting, nor no evil speaking; and see that the church meets together oft, and also see that every member does his duty, and he is to take the lead of the meetings in the absence of the elder or priest, and is to be assisted always and all his duties in the church by the deacons. But neither the teacher nor the deacon has authority to baptize nor administer the sacrament; but are to warn, exhort, expound and teach and invite all to come to Christ. Every elder, priest, teacher, or deacon, is to be ordained according to the gifts and calling of God unto them by the power of the Holy Ghost, which is in the one who ordains them. The several elders composing the church of Christ are to meet at each of its meetings to do church business, whatsoever is necessary, \&c. and each priest or teacher who is ordained by any priest is to take a certificate from him at the time, which, when shown to an elder, he is to give him a license, which shall authorize him to perform the duty of his calling. The duty of the church members after they are received by baptism:---The elders or priests are to have a sufficient time to expound all things concerning this church of Christ to their understanding previous to their partaking of the sacrament, and being confirmed by the laying on of the hands of the elders, so that all things shall be done in order; and the members shall manifest before the church and before the elders a godly walk and conversation that they are worthy of it, that there may be works and faith agreeable to the holy scriptures, walking in holiness before the Lord. Every member of this church of Christ having children, are to bring them unto the elders before the church who are to lay hands on them in the name of the Lord, and bless them in the name of Christ. There cannot any one be received into this church of Christ who have not arrived to the years of accountability before God, and are not capable of repentance. And the manner of baptism \& the manner of administering the sacrament are to be done as is written in the Book of Morman. Any member of this church of Christ transgressing, or being overtaken in a fault, shall be dealt with according as the scriptures direct, \&c It shall be the duty of the several churches composing this church of Christ to send one of their priests or teachers to attend the several conferences held by the elders of the church with a list of the names the several persons uniting themselves to the church since the last conference, or send by the hand of some elder, so that there can be kept a regular list of all the names of the members of the whole church in a book kept by one of the elders whomsoever the other elders shall appoint from time to time, end [and] also if any have been expelled from the church so that their names may be blotted out of the general church record of names; any member removing from the church where he belongs, if going to a church where he is not known, may take letter certifying that he is a member and in good standing, which certificate may be signed by any elder or priest---if the person receiving the letter is personally acquainted with the elder or priest; or may be signed by the teachers or deacons of his church.

		% -----
		\subsubsection{Revelation, 6 April 1830 [D\&C 21]}
		% -----
			\label{EMS-1830-JS-6-APR-1830}
			% \label{EMS-1830-JS-DAY-3LETTERMONTH-YEAR}
		
			\begin{description}
				\item[Print Sources]
			\end{description}

				\begin{tabular}{p{0.5in}p{1in}p{0.5in}p{1in}}
				JSP	 	& D1:126--130	& JSR 		& 60--61	 \\
				DC	 	& Section 21	& HC 		& 1:78--79	 \\
				HDDC 	& 352--359		& TCHC 		& 1:56 \\
				\end{tabular}
				% HDDC = woodford's DC diss; JSR = Marquardt's JS Revelations; TCHC = Vogel HC
			
			\begin{description}
				\item[Online Access] \href{https://www.josephsmithpapers.org/paper-summary/revelation-6-april-1830-dc-21/1}{Here}
				% \item[Analysis] \hyperref[DC21]{Here}
			\end{description}
			
			\begin{description}
				\item[Background]
			\end{description}
			
				% It might also be useful to give a two sentence context of the period: e.g., "This speech was given in Nauvoo to a small congregation one month prior to the destruction of the Nauvoo Expositor. These and other tensions may have been on the prophet's mind when he gave the speech."
			
			\begin{description}
				\item[Original Source]
			\end{description}
			
				% brief description of original source material, with reference to JSP or somewhere else for more info
				% Background material: it might be helpful to have a note that says whether a given document was presented as a speech, was written in a journal, etc.
				% Also a comment here about how reliable the source might be, or what shape it is in, if there are large omissions or parts that are hard to understand, and so on.
			
			\begin{description}
				\item[Text]
			\end{description}
			
				\begin{tightcenter}
				17\supersu{th}\supers{.} Commandment AD 1829
				\end{tightcenter}
				
				A Revelation to Joseph the Seer by way of commandment to the Church given at Fayette Seneca County State of New York
				
				Behold there Shall a Record be kept among you \& in it thou shalt be called a seer \& Translater \& Prop[h]et an Apostle of Jesus Christ an Elder of the Church through the will of God the Father \& the grace of our Lord Jesus Christ being inspired of the Holy Ghost to lay the foundation thereof \& to build it up unto the most holy faith which Church was Organized \& established in the year of our Lord one thousand Eight Hundred \& Thirty in the forth Month on the Sixth day of the month which is called April Wherefore meaning the Church thou shalt give heed unto all his words \& commandments which he ghall [shall] give unto you as he receiveth them wa[l]king in all holyness before me for his word ye shall receive as if from mine own mouth in all Patience \& faith for by doing these things the gaits of Hell shall not prevail against you yea \& the Lord God will disperse the Powers of darkness from before you \& cause the Heavens to shake for your Good \& his names glory for thus saith the Lord God him have I inspired to move the cause of Zion in Mighty power for good \& his dilligence I know \& his prayers I have heard yea his weeping for Zion I have seen \& I will cause that He shall mourn for her no longer for his days of rejoicing are come unto the remission of his Sins \& the manifestations of my blessings upon his works for behold I will bless all those who Labour in my Vinyard with a mighty blessing \& they shall believe on his words which are given him through me by the comforter which manifesteth that Jesus was Crusified by the Sins of the world for the remision of sins unto the contrite heart Wherefore it behooveth me that he should be ordained by you Oliver [Cowdery] mine Apostle this being an Ordinance unto you that ye are an Elder under his hand unto you that thou mightest be an Elder unto this Church of Christ bearing my name \& the first Preacher of this Church unto the Church \& before the world yea before the gentiles yea \sout{before} \& thus saith the Lord God Lo. Lo. to the Jews also Amen---
				
		% -----
		\subsubsection{Revelation, April 1830--A [D\&C 23:1--2]}
		% -----
			\label{EMS-1830-JS-APR-1830-A}
			% \label{EMS-1830-JS-DAY-3LETTERMONTH-YEAR}
		
			\begin{description}
				\item[Print Sources]
			\end{description}

				\begin{tabular}{p{0.5in}p{1in}p{0.5in}p{1in}}
				JSP	 	& D1:130--131	& JSR 		& 57--59	 \\
				DC	 	& Section 23	& HC 		& 1:80		 \\
				HDDC 	& 367--372		& TCHC 		& 1:57--58 \\
				\end{tabular}
				% HDDC = woodford's DC diss; JSR = Marquardt's JS Revelations; TCHC = Vogel HC
			
			\begin{description}
				\item[Online Access] \href{https://www.josephsmithpapers.org/paper-summary/revelation-april-1830-a-dc-231-2/1}{Here}
				% \item[Analysis] \hyperref[DC23]{Here}
			\end{description}
			
			\begin{description}
				\item[Background]
			\end{description}
			
				% It might also be useful to give a two sentence context of the period: e.g., "This speech was given in Nauvoo to a small congregation one month prior to the destruction of the Nauvoo Expositor. These and other tensions may have been on the prophet's mind when he gave the speech."
			
			\begin{description}
				\item[Original Source]
			\end{description}
			
				% brief description of original source material, with reference to JSP or somewhere else for more info
				% Background material: it might be helpful to have a note that says whether a given document was presented as a speech, was written in a journal, etc.
				% Also a comment here about how reliable the source might be, or what shape it is in, if there are large omissions or parts that are hard to understand, and so on.
			
			\begin{description}
				\item[Text]
			\end{description}
			
				\begin{tightcenter}
				18\supersu{th}\supers{.} Commandment AD 1830
				\end{tightcenter}
				
				A Revelation to Oliver [Cowdery] given at Manchester Ontario Co State of New York Soon after his calling to the Minstery

				Behold I say unto you Oliver a few words Behold thou art Blessed \& art under no condemnation but beware of pride lest thou shouldest enter into temptation make known thy Calling unto the Church \& also before the World \& thy heart shall be opened to Preach the truth from henceforth \& for ever amen
				
		% -----
		\subsubsection{Revelation, April 1830--B [D\&C 23:3]}
		% -----
			\label{EMS-1830-JS-APR-1830-B}
			% \label{EMS-1830-JS-DAY-3LETTERMONTH-YEAR}
		
			\begin{description}
				\item[Print Sources]
			\end{description}

				\begin{tabular}{p{0.5in}p{1in}p{0.5in}p{1in}}
				JSP	 	& D1:132		& JSR 		& 59	 \\
				DC	 	& Section 23	& HC 		& 1:80	 \\
				HDDC 	& 367--372		& TCHC 		& 1:57--58 \\
				\end{tabular}
				% HDDC = woodford's DC diss; JSR = Marquardt's JS Revelations; TCHC = Vogel HC
			
			\begin{description}
				\item[Online Access] \href{https://www.josephsmithpapers.org/paper-summary/revelation-april-1830-b-dc-233/1}{Here}
				% \item[Analysis] \hyperref[DC23]{Here}
			\end{description}
			
			\begin{description}
				\item[Background]
			\end{description}
			
				% It might also be useful to give a two sentence context of the period: e.g., "This speech was given in Nauvoo to a small congregation one month prior to the destruction of the Nauvoo Expositor. These and other tensions may have been on the prophet's mind when he gave the speech."
			
			\begin{description}
				\item[Original Source]
			\end{description}
			
				% brief description of original source material, with reference to JSP or somewhere else for more info
				% Background material: it might be helpful to have a note that says whether a given document was presented as a speech, was written in a journal, etc.
				% Also a comment here about how reliable the source might be, or what shape it is in, if there are large omissions or parts that are hard to understand, and so on.
			
			\begin{description}
				\item[Text]
			\end{description}
			
				\begin{tightcenter}
				19\supersu{th}\supers{.} Commandment AD 1830
				\end{tightcenter}
				
				A Commandment to Hyram [Hyrum Smith] given at Manchester Ontario County State of New York
				
				Behold I speak unto you Hyram a few words for thou also art under no condemnation \& thy heart is opened \& thy tongue loosed \& thy Calling is to exhortation \& to strengthen the Church continually wherefore thy duty is unto the Church forever \& this because of thy family
				
		% -----
		\subsubsection{Revelation, April 1830--C [D\&C 23:4]}
		% -----
			\label{EMS-1830-JS-APR-1830-C}
			% \label{EMS-1830-JS-DAY-3LETTERMONTH-YEAR}
		
			\begin{description}
				\item[Print Sources]
			\end{description}

				\begin{tabular}{p{0.5in}p{1in}p{0.5in}p{1in}}
				JSP	 	& D1:133		& JSR 		& 59	 \\
				DC	 	& Section 23	& HC 		& 1:80	 \\
				HDDC 	& 367--372		& TCHC 		& 1:57--58 \\
				\end{tabular}
				% HDDC = woodford's DC diss; JSR = Marquardt's JS Revelations; TCHC = Vogel HC
			
			\begin{description}
				\item[Online Access] \href{https://www.josephsmithpapers.org/paper-summary/revelation-april-1830-c-dc-234/1}{Here}
				% \item[Analysis] \hyperref[DC23]{Here}
			\end{description}
			
			\begin{description}
				\item[Background]
			\end{description}
			
				% It might also be useful to give a two sentence context of the period: e.g., "This speech was given in Nauvoo to a small congregation one month prior to the destruction of the Nauvoo Expositor. These and other tensions may have been on the prophet's mind when he gave the speech."
			
			\begin{description}
				\item[Original Source]
			\end{description}
			
				% brief description of original source material, with reference to JSP or somewhere else for more info
				% Background material: it might be helpful to have a note that says whether a given document was presented as a speech, was written in a journal, etc.
				% Also a comment here about how reliable the source might be, or what shape it is in, if there are large omissions or parts that are hard to understand, and so on.
			
			\begin{description}
				\item[Text]
			\end{description}
			
				\begin{tightcenter}
				\sout{29} 20\supers{th} Commandment AD 1830
				\end{tightcenter}
				
				A Revelation to Samuel [Smith] given at Manchester Ontario Co N.Y.
				
				Behold I speak a few words unto you Samuel for thou also art under no condemnation \& thy calling is to Exhortation to strengthen the Church \& thou art not <as> yet called to prea[c]h before the world Amen
				
		% -----
		\subsubsection{Revelation, April 1830--D [D\&C 23:5]}
		% -----
			\label{EMS-1830-JS-APR-1830-D}
			% \label{EMS-1830-JS-DAY-3LETTERMONTH-YEAR}
		
			\begin{description}
				\item[Print Sources]
			\end{description}

				\begin{tabular}{p{0.5in}p{1in}p{0.5in}p{1in}}
				JSP	 	& D1:133--134	& JSR 		& 60	 \\
				DC	 	& Section 23	& HC 		& 1:80	 \\
				HDDC 	& 367--372		& TCHC 		& 1:57--58 \\
				\end{tabular}
				% HDDC = woodford's DC diss; JSR = Marquardt's JS Revelations; TCHC = Vogel HC
			
			\begin{description}
				\item[Online Access] \href{https://www.josephsmithpapers.org/paper-summary/revelation-april-1830-d-dc-235/1}{Here}
				% \item[Analysis] \hyperref[DC23]{Here}
			\end{description}
			
			\begin{description}
				\item[Background]
			\end{description}
			
				% It might also be useful to give a two sentence context of the period: e.g., "This speech was given in Nauvoo to a small congregation one month prior to the destruction of the Nauvoo Expositor. These and other tensions may have been on the prophet's mind when he gave the speech."
			
			\begin{description}
				\item[Original Source]
			\end{description}
			
				% brief description of original source material, with reference to JSP or somewhere else for more info
				% Background material: it might be helpful to have a note that says whether a given document was presented as a speech, was written in a journal, etc.
				% Also a comment here about how reliable the source might be, or what shape it is in, if there are large omissions or parts that are hard to understand, and so on.
			
			\begin{description}
				\item[Text]
			\end{description}
			
				\begin{tightcenter}
				21\supers{st} Commandment AD 1830
				\end{tightcenter}
				
				A Commandment to given to Joseph [Smith Sr.] at Manchester Ontario County State of New York
				
				Behold I speak a few words \sout{unto you Joseph f} unto you Joseph for thou art under no condemnation \& thy calling also is to Exhortation \& to strengthen the Church \& this is thy duty from henceforth \& forever amen
				
		% -----
		\subsubsection{Revelation, April 1830--E [D\&C 23:6--7]}
		% -----
			\label{EMS-1830-JS-APR-1830-E}
			% \label{EMS-1830-JS-DAY-3LETTERMONTH-YEAR}
		
			\begin{description}
				\item[Print Sources]
			\end{description}

				\begin{tabular}{p{0.5in}p{1in}p{0.5in}p{1in}}
				JSP	 	& D1:134--136	& JSR 		& 60	 \\
				DC	 	& Section 23	& HC 		& 1:80	 \\
				HDDC 	& 367--372		& TCHC 		& 1:57--58 \\
				\end{tabular}
				% HDDC = woodford's DC diss; JSR = Marquardt's JS Revelations; TCHC = Vogel HC
			
			\begin{description}
				\item[Online Access] \href{https://www.josephsmithpapers.org/paper-summary/revelation-april-1830-e-dc-236-7/1}{Here}
				% \item[Analysis] \hyperref[DC23]{Here}
			\end{description}
			
			\begin{description}
				\item[Background]
			\end{description}
			
				% It might also be useful to give a two sentence context of the period: e.g., "This speech was given in Nauvoo to a small congregation one month prior to the destruction of the Nauvoo Expositor. These and other tensions may have been on the prophet's mind when he gave the speech."
			
			\begin{description}
				\item[Original Source]
			\end{description}
			
				% brief description of original source material, with reference to JSP or somewhere else for more info
				% Background material: it might be helpful to have a note that says whether a given document was presented as a speech, was written in a journal, etc.
				% Also a comment here about how reliable the source might be, or what shape it is in, if there are large omissions or parts that are hard to understand, and so on.
			
			\begin{description}
				\item[Text]
			\end{description}
			
				\begin{tightcenter}
				22\supers{nd} Commandment AD 1830
				\end{tightcenter}
				
				A Commandment to Joseph Knight [Sr.] given at Manchester Ontario County State of New York
				
				Behold I manifest unto you <by> these words that thou must take up thy Cross in the which thou must pray vocally before the World as well as in Seecret \& in thy family \& among thy friends \& in all Places Behold it is thy duty to unite with the true Church \& give thy Language to Exhortation continually that thou mayest Receive the reward of the Labourer amen
				
		% -----
		\subsubsection{Revelation, 16 April 1830 [D\&C 22]}
		% -----
			\label{EMS-1830-JS-16-APR-1830}
			% \label{EMS-1830-JS-DAY-3LETTERMONTH-YEAR}
		
			\begin{description}
				\item[Print Sources]
			\end{description}

				\begin{tabular}{p{0.5in}p{1in}p{0.5in}p{1in}}
				JSP	 	& D1:137--138	& JSR 		& 62	 	\\
				DC	 	& Section 22	& HC 		& 1:79--80 	\\
				HDDC 	& 360--366		& TCHC 		& 1:57 		\\
				\end{tabular}
				% HDDC = woodford's DC diss; JSR = Marquardt's JS Revelations; TCHC = Vogel HC
			
			\begin{description}
				\item[Online Access] \href{https://www.josephsmithpapers.org/paper-summary/revelation-16-april-1830-dc-22/1}{Here}
				% \item[Analysis] \hyperref[DC22]{Here}
			\end{description}
			
			\begin{description}
				\item[Background]
			\end{description}
			
				% It might also be useful to give a two sentence context of the period: e.g., "This speech was given in Nauvoo to a small congregation one month prior to the destruction of the Nauvoo Expositor. These and other tensions may have been on the prophet's mind when he gave the speech."
			
			\begin{description}
				\item[Original Source]
			\end{description}
			
				% brief description of original source material, with reference to JSP or somewhere else for more info
				% Background material: it might be helpful to have a note that says whether a given document was presented as a speech, was written in a journal, etc.
				% Also a comment here about how reliable the source might be, or what shape it is in, if there are large omissions or parts that are hard to understand, and so on.
			
			\begin{description}
				\item[Text]
			\end{description}
			
				A commandment unto the church of Christ which was established in these last days A. D. 1830, on the 4th month and the 6th day of the month, which is called April:--- Behold I say unto you that all old covenants have I caused to be done away in this thing, and this is a new and everlasting covenant; wherefore, although a man should be baptized a hundred times it availeth him nothing, for ye cannot enter in at the straight gait by the law of Moses, neither by your dead works, for it is because of your dead works that I have caused this last covenant and this church to be built up unto me; wherefore enter ye in at the gate as I have commanded, and seek not to counsel your God.
				
		% -----
		\subsubsection{Visions of Moses, June 1830 [Moses 1]}
		% -----
			\label{EMS-1830-JS-JUN-1830}
			% \label{EMS-1830-JS-DAY-3LETTERMONTH-YEAR}
		
			\begin{description}
				\item[Print Sources]
			\end{description}

				\begin{tabular}{p{0.5in}p{1in}p{0.5in}p{1in}}
				JSP	 	& D1:150--156	& JSR 		& --- 		\\
				DC	 	& ---			& HC 		& 1:98--101 \\
				HDDC 	& ---			& TCHC 		& 1:70--71 \\
				\end{tabular}
				% HDDC = woodford's DC diss; JSR = Marquardt's JS Revelations; TCHC = Vogel HC
			
			\begin{description}
				\item[Online Access] \href{https://www.josephsmithpapers.org/paper-summary/visions-of-moses-june-1830-moses-1/1}{Here}
				% \item[Analysis] \hyperref[XX]{Here}
			\end{description}
			
			\begin{description}
				\item[Background]
			\end{description}
			
				% It might also be useful to give a two sentence context of the period: e.g., "This speech was given in Nauvoo to a small congregation one month prior to the destruction of the Nauvoo Expositor. These and other tensions may have been on the prophet's mind when he gave the speech."
			
			\begin{description}
				\item[Original Source]
			\end{description}
			
				% brief description of original source material, with reference to JSP or somewhere else for more info
				% Background material: it might be helpful to have a note that says whether a given document was presented as a speech, was written in a journal, etc.
				% Also a comment here about how reliable the source might be, or what shape it is in, if there are large omissions or parts that are hard to understand, and so on.
			
			\begin{description}
				\item[Text]
			\end{description}
			
				\begin{tightcenter}
				A Revelation given to Joseph the Revelator June 1830
				\end{tightcenter}
				
				The words of God which he \sout{gave} <spake> unto Moses at a time when Moses was caught up into an exceeding high Mountain \& he saw God face to face \& he talked with him \& the glory of God was upon Moses therefore Moses could endure his presence \& God spake unto Moses saying Behold I I am the Lord God Almighty \& endless is my name for I am without beginning of days or end of years \& is this not endless \& behold thou art my Son Wherefore look \& I will shew thee the workmanship of mine hands but not all for my works are without end \& also my wo[r]ds for they never cease wherefore no man can behold all my works except he behold <all> my glory \& no man can behold all my glory \& afterwords remain in the flesh \& I have a work for thee Moses my Son \& thou art in similitude to \sout{my} <mine> only begotten \& mine only begotten is \& shall be for he is full of grace \& truth but there is none other God beside me \& all things are present with me for I know them all And now behold this one thing I shew unto thee Moses my son for thou art in the world \& now I shew it thee And it came to pass that Moses looked \& beheld the world upon which he was created \& Moses beheld the world \& the ends thereof \& all the Children of men which was \& which was created of the same he greatly marvelled \& wondered \& the presence of God withdrew from Moses that his glory was not upon Moses \& Moses was left unto himself \& as he was left unto himself he fell unto the Earth, And it came to pass, that it was for the space of many hours before Moses did again receive his natural strength [lik]e unto man, \& he saith unto himself \sout{no} Now for this once I know that [m]an is nothing, which thing I never had supposed, but now mine eyes, mine own eyes but not mine eyes for mine eyes could not have beheld for I should have withered \& died in his presence but his glory was upon me \& I beheld his face for I was transfigered before him And now it came to pass that when Moses had said these words, behold, Satan came tempting him saying, Moses, Son of man, worship me, And it came to pass that Moses looked upon Satan \& saith, Who art thou for behold, I am a Son of God in the similitude of his only begotten, \& where is thy glory that I should worship thee, for, behold, I could not look upon God except his glory should come upon me, \& I were transfigered before him but I can look upon thee in the natural man, if not so surely blessed be the name of my God for his Spirit hath not altogether withdrawn from me or else where is thy glory for it is blackness unto me \& I can Judge between thee \& God for God said unto me Worship God for him only shalt thou serve Get thee hence Satan deceive me not for God said unto me Thou art after the similitude of mine only begotten \& he also gave unto me commandment when he called unto me out of the burning bush Saying [c]all upon God in the name of mine only begotten \& worship me And again Moses saith I will not cease to call upon God I have other things to inquire of him for his glory has been upon me \& it is glory unto me wherefore I can judge betwixt him \& thee[.] depart hence Satan And now when Mose[s] had said these words Satan cried with a loud [voice \&] wrent upon the Ea[r]t[h] \& comma[n]ded saying I am the only begotten worship me And it came to pass that Moses began to fear exceedingly \& as he began to fear he saw the bitterness of Hell Nevertheless calling upon God he received strength \& he commanded saying Depart hence Satan for this one God only will I worship which is the God of glory \& now Satan began to tremble \& the Earth shook \& Moses receiving strength called upon God saying In the name of Jesus Christ depart hence Satan And it came to pass that satan cried with a loud voice with weeping \& wailing \& gnashing of teeth \& departed hence yea from the presence of Moses that he beheld him not And now of this \sout{thing} <thing> Moses bore record but because of wickedness it is not had among the children of men And it came to pass that when Satan had departed from the presence of Moses he lifted up his eyes unto Heaven being filled with the Holy Ghost which beareth record of the Father \& the Son \& calling upon the name of God he beheld again his glory for it was upon him \& he heard a voice saying Blessed art thou Moses for I the Almighty have chosen thee \& thou shalt be made stronger than the many waters for they shall obey thy command even as if thou wert God \& lo I am with you even to the end of thy days for thou shalt deliver my people from bondage even Israel my chose[n] And it came to pass as the voice was still speaking he cast his eyes \& h[e] beheld the Earth yea even all all the face of it \& there was not a particl[e] of it which he did not behold diserning it by the Spirit of God \& he beheld also the inhabitants thereof \& there was not a soul which he be[held] not \& he discerned them by the Spirit of God \& their numbers were grea[t] even as numberless as the sand upon the sea shore \& he beheld many lands \& each land was called Earth \& there were inhabitants upon the face thereof And it came to pass that Moses called upon God saying tell me I pray thee why these things are so \& by what thou madest them \& behold the glory of God was upon Moses that Moses stood in the presence of God \& he talk[ed] with him face to face \& the Lord God said unto Moses For mine own purpose have I made these things here is wisdom \& it remaineth in me \& by the word of my power have I created them which is mine only begotten Son full of grace \& truth \& worlds without number have I created \& I also created them for mine own purpose \& by the same I created them which is mine only begotten \& the first man of all men have I called Adam which is many but only an account of this Earth \& the inhabitants thereof give I unto you for behold there are many worlds which have passed away by the word of my power \& there are many also which now stand \& numberless are they unto man bu[t] all things are numbered unto me for they are mine \& I know them And it came to pass that Moses spake unto the Lord saying Be merciful [u]nto thy Servant O God \& tell me concerning this Earth \& the inhabitan[ts] [the]reof \& also the H[eavens] \& then thy Servant will be content And the Lord God spake unto Moses saying The Heavens there are many \& they cannot be numbered unto man but they are numbered unto me for they are mine \& as one Earth shall pass away \& the Heavens thereof even so shall another come And there is no end to my works neither my words for behold this is my work to my glory to the immortality \& the eternal life of man And now Moses my Son I will speak unto you concerning this Earth upon which thou standest \& thou shalt write the things which I shall speak \& in a day when the children of men shall esteem my words as nought \& take many of them from the Book which thou shalt write behold I will raise up another like unto \sout{theee} thee \& they shall be had again among the Children of men among even as many as shall believe These words was spoken unto Moses in the mount the name of which shall not be known among the Children of men And now they are also spoken unto you shew them not unto any except them that believe <Amen>
				
		% -----
		\subsubsection{Revelation, July 1830--A [D\&C 24]}
		% -----
			\label{EMS-1830-JS-JUL-1830-A}
			% \label{EMS-1830-JS-DAY-3LETTERMONTH-YEAR}
		
			\begin{description}
				\item[Print Sources]
			\end{description}

				\begin{tabular}{p{0.5in}p{1in}p{0.5in}p{1in}}
				JSP	 	& D1:156--159	& JSR 		& 69--70		\\
				DC	 	& Section 24	& HC 		& 1:101--103 	\\
				HDDC 	& 373--380		& TCHC 		& 1:72--73 			\\
				\end{tabular}
				% HDDC = woodford's DC diss; JSR = Marquardt's JS Revelations; TCHC = Vogel HC
			
			\begin{description}
				\item[Online Access] \href{https://www.josephsmithpapers.org/paper-summary/revelation-july-1830-a-dc-24/1}{Here}
				% \item[Analysis] \hyperref[DC24]{Here}
			\end{description}
			
			\begin{description}
				\item[Background]
			\end{description}
			
				% It might also be useful to give a two sentence context of the period: e.g., "This speech was given in Nauvoo to a small congregation one month prior to the destruction of the Nauvoo Expositor. These and other tensions may have been on the prophet's mind when he gave the speech."
			
			\begin{description}
				\item[Original Source]
			\end{description}
			
				% brief description of original source material, with reference to JSP or somewhere else for more info
				% Background material: it might be helpful to have a note that says whether a given document was presented as a speech, was written in a journal, etc.
				% Also a comment here about how reliable the source might be, or what shape it is in, if there are large omissions or parts that are hard to understand, and so on.
			
			\begin{description}
				\item[Text]
			\end{description}
			
				\begin{tightcenter}
				25\supersu{th}\supers{.} Commandment AD \sout{1831} 1830
				\end{tightcenter}
				
				A Revelation to Joseph \& Oliver [Cowdery] given at Harmony Susquehannah County Pennsylvania telling them concerning their Calls \&c
				
				Behold thou wast called \& Chosen to write the Book of Mormon \& to my ministery \& I have lifted thee up out of thine afflictions \& have counseled thee that thou hast been delivered from all thine enemies \& thou hast been delivered from the power of satan \& from darkness Nevertheless thou art not excusable in thy Transgressions Nevertheless go thy way \& sin no more magnify thy office \& after that thou hast sowed thy fields \& Secured them then go speedily unto the Church which is in Colesvill Fayette \& Manchester \& they shall support thee \& I will bless them both spiritually \& temporally but if they receive thee not I will send upon them a cursing instead of a blessing \& thou shalt continue in calling upon [God] in my name \& writing the Things which shall be given thee by the Comforter \& thou shalt expound all scriptures \sout{to} unto the Church \& it shall be given thee in the very moment what thou shalt speak \& write \& they shall hear it or I will send unto them a cursing instead of a blessing for thou shalt devote all thy service \sout{to} in Zion \& in this thou shalt have strength be patient in afflictions for thou shalt have many but endure them for Lo! I am with thee even unto the end of thy days \& in temporal labo$\diamond\diamond$ thou shalt not have strength for this is not thy calling attend to thy calling \& thou shalt have wherewith to magnify thine Office \& to expound all scriptures \& continue in the laying \sout{of} on of the hands \& confirming the Churches \& thy brother Oliver Shall continue \sout{to} in bearing my name before the world \& also to the Church \& he shall not suppose that he can say enough in my cause \& lo! [I] am with him to the end in me he shall have glory \& not of himself whether in weakness or in strength whether in bonds \sout{on} or free \& at all times \& in all places he shall open his mouth \& declare my Gospel as with the voice of a Trump both day \& night \& I will give unto him strength such as is not known among men[.] require not Miracles except I shall command you except casting [out] Devils healing the sick \& against Poisones [poisonous] Serpents \& against deadly Poison \& these things ye shall not do except it be required of you by them who desire it that the Scriptures might be fulfilled for ye shall do according to that which is written in the Scriptures \& in whatsoever place ye shall enter in \& they receive you not in my name ye shall leave a cursing instead of a blessing by casting off the dust of your feet against them as a testimony \& cleansing your feet by the wayside \& it shall come to pass that whosoever shall lay their hands upon you by violence ye shall command to be smitten <in my name> \& behold I will smite them according to \sout{their} <thy> words in mine own due time \& whosoever shall go to law with thee shall be cursed by the Law \& thou shalt tak[e] no purse nor scrip neither staves neither two Coats for the Church shall give unto thee in the verry hour what thou needest for food \& for raiment for shoes \& for Money \& for scrip for thou art called to prune my vineyard with a mighty pruneing yea even for the last time yea \& also all those \sout{which} whom thou hast ordained \& they shall do even according to this pattern amen
				
		% -----
		\subsubsection{Revelation, July 1830--B [D\&C 26]}
		% -----
			\label{EMS-1830-JS-JUL-1830-B}
			% \label{EMS-1830-JS-DAY-3LETTERMONTH-YEAR}
		
			\begin{description}
				\item[Print Sources]
			\end{description}

				\begin{tabular}{p{0.5in}p{1in}p{0.5in}p{1in}}
				JSP	 	& D1:160--161	& JSR 		& 71	 \\
				DC	 	& Section 26	& HC 		& 1:104	 \\
				HDDC 	& 387--392		& TCHC 		& 1:73--74 \\
				\end{tabular}
				% HDDC = woodford's DC diss; JSR = Marquardt's JS Revelations; TCHC = Vogel HC
			
			\begin{description}
				\item[Online Access] \href{https://www.josephsmithpapers.org/paper-summary/revelation-july-1830-b-dc-26/1}{Here}
				% \item[Analysis] \hyperref[DC26]{Here}
			\end{description}
			
			\begin{description}
				\item[Background]
			\end{description}
			
				% It might also be useful to give a two sentence context of the period: e.g., "This speech was given in Nauvoo to a small congregation one month prior to the destruction of the Nauvoo Expositor. These and other tensions may have been on the prophet's mind when he gave the speech."
			
			\begin{description}
				\item[Original Source]
			\end{description}
			
				% brief description of original source material, with reference to JSP or somewhere else for more info
				% Background material: it might be helpful to have a note that says whether a given document was presented as a speech, was written in a journal, etc.
				% Also a comment here about how reliable the source might be, or what shape it is in, if there are large omissions or parts that are hard to understand, and so on.
			
			\begin{description}
				\item[Text]
			\end{description}
			
				\begin{tightcenter}
				26\supers{th} Commandment AD 1830
				\end{tightcenter}
				
				A Revelation to Joseph Oliver [Cowdery] \& John [Whitmer] given at Harmony Susquehannah County State of Pennsylvania
				
				Behold I say unto you that ye shall let your time be devoted to the studying the Scriptures \& to preaching \& to confirming the Church at Colesvill \& to performing thy labours on the Land such as is required until after ye shall go to the west to hold the next conference then it shall be made known what thou shalt do \& all things shall be done by common consent in the Church by much prayer \& faith for all things ye shall receive by faith \& thus it is amen
				
		% -----
		\subsubsection{Revelation, July 1830--C [D\&C 25]}
		% -----
			\label{EMS-1830-JS-JUL-1830-C}
			% \label{EMS-1830-JS-DAY-3LETTERMONTH-YEAR}
		
			\begin{description}
				\item[Print Sources]
			\end{description}

				\begin{tabular}{p{0.5in}p{1in}p{0.5in}p{1in}}
				JSP	 	& D1:161--164	& JSR 		& 70--71		\\
				DC	 	& Section 25	& HC 		& 1:103--104 	\\
				HDDC 	& 381--386		& TCHC 		& 1:73 			\\
				\end{tabular}
				% HDDC = woodford's DC diss; JSR = Marquardt's JS Revelations; TCHC = Vogel HC
			
			\begin{description}
				\item[Online Access] \href{https://www.josephsmithpapers.org/paper-summary/revelation-july-1830-c-dc-25/1}{Here}
				% \item[Analysis] \hyperref[DC25]{Here}
			\end{description}
			
			\begin{description}
				\item[Background]
			\end{description}
			
				% It might also be useful to give a two sentence context of the period: e.g., "This speech was given in Nauvoo to a small congregation one month prior to the destruction of the Nauvoo Expositor. These and other tensions may have been on the prophet's mind when he gave the speech."
			
			\begin{description}
				\item[Original Source]
			\end{description}
			
				% brief description of original source material, with reference to JSP or somewhere else for more info
				% Background material: it might be helpful to have a note that says whether a given document was presented as a speech, was written in a journal, etc.
				% Also a comment here about how reliable the source might be, or what shape it is in, if there are large omissions or parts that are hard to understand, and so on.
			
			\begin{description}
				\item[Text]
			\end{description}
				
				\begin{tightcenter}
				27\supersu{th}\supers{.} Commandment AD 1830
				\end{tightcenter}
				
				A Revelation to Emma [Smith] given at Harmony Susquehan[na] County state of Pennsylvania giving her a command to select Hymns \&c
				
				A Revelation I give unto you concerning my will Behold thy sins are for given thee \& thou art an Elect Lady whom I have called murmer not because of the things which thou hast not seen for they are withheld from thee \& the World which is wisdom in me in a time to come \& the office of thy calling shall be for a comfort unto my Servent Joseph thy husband in his afflictions with consoleing words in the spirit of meekness \& thou shalt go with him at the time of his going \& be unto him a Scribe that I may send Oliver [Cowdery] whithersoever I will \& thou shalt be ordained under his hand to expound Scriptures \& exhort the Church according as it shall be given thee by my spirit for he shall lay his hands upon the[e] \& thou shalt receive the Holy Ghost \& thy time shall be Given to writing\sout{s} \& to Learning \& thou needest not fear for thy husband shall support thee from the Church for unto them is \sout{thy} his calling that all things might be revealed unto them whatsoever I will according to their faith \& verily I say unto \sout{you} thee that thou shalt lay aside the things of this world \& seek for the things of a better \& it shall be given thee also to make a selection of Sacred Hymns as it shall be given thee which is pleasing unto me to be had in my Church for my Soul delighteth in the song of the heart yea the song of the \sout{heart} righteous is a prayer unto me \& it shall be answered with a blessing upon their heads wherefore lift up thy heart \& rejoice \& cleave unto the covenants which thou hast made continue in the spirit of meekness \& beware of Pride let thy soul delight in thy husband \& the glory which shall come upon him keep my commandments continually \& a crown of righteousness thou shalt receive \& except thou do this where I am \sout{thou} ye cannot come \& verily I say unto you that this is my voice unto all even so amen
				
		% -----
		\subsubsection{Revelation, circa August 1830 [D\&C 27]}
		% -----
			\label{EMS-1830-JS-CIR-AUG-1830}
			% \label{EMS-1830-JS-DAY-3LETTERMONTH-YEAR}
		
			\begin{description}
				\item[Print Sources]
			\end{description}

				\begin{tabular}{p{0.5in}p{1in}p{0.5in}p{1in}}
				JSP	 	& D1:164--166	& JSR 		& 72--80\footnote{Note about how Marquardt includes some extensive commentary here.} \\
				DC	 	& Section 27	& HC 		& 1:106--108 \\
				HDDC 	& 393--403		& TCHC 		& 1:75--76 \\
				\end{tabular}
				% HDDC = woodford's DC diss; JSR = Marquardt's JS Revelations; TCHC = Vogel HC
			
			\begin{description}
				\item[Online Access] \href{https://www.josephsmithpapers.org/paper-summary/revelation-circa-august-1830-dc-27/1}{Here}
				% \item[Analysis] \hyperref[DC27]{Here}
			\end{description}
			
			\begin{description}
				\item[Background]
			\end{description}
			
				% It might also be useful to give a two sentence context of the period: e.g., "This speech was given in Nauvoo to a small congregation one month prior to the destruction of the Nauvoo Expositor. These and other tensions may have been on the prophet's mind when he gave the speech."
			
			\begin{description}
				\item[Original Source]
			\end{description}
			
				% brief description of original source material, with reference to JSP or somewhere else for more info
				% Background material: it might be helpful to have a note that says whether a given document was presented as a speech, was written in a journal, etc.
				% Also a comment here about how reliable the source might be, or what shape it is in, if there are large omissions or parts that are hard to understand, and so on.
			
			\begin{description}
				\item[Text]
			\end{description}
			
				\begin{center}
				28\supersu{th}\supers{.} Commandment AD 1830
				\end{center}
				
				A Revelation to the Church given at Harmony susquehannh County State of Pennsylvania given to Joseph the Seer at a time that he went to purchase wine \sout{it} for Sacrament \& he was stoped by an Angel \& \sout{he} he spok to him as follows Saying
				
				Listen to the voice of Jesus Christ your Lord your God \& your Redeemer whose word is quick \& powerful for Behold I say unto you it mattereth not what ye <shall> eat or what ye shall drink when ye partake of the sacrament if it so be that ye do it with an eye single to my glory Remembering unto the father my Body which [was] laid down for you \& my blood which was shed for \sout{you} the Remission of your sins Wherefore a commandment I give unto you that ye shall not Purchase Wine neither strong drink of your enemies Wherefore ye shall partake none except it is made new among you yea in this my Fathers Kingdom which shall be built up on the earth Behold this is wisdom in me Wherefore marvel not for the hour cometh that I will drink of the fruit of the Vine with you on the Earth \& with \sout{you} all those whom my father hath given me out of the world Wherefore lift up your hearts \& rejoice \& Gird up your loins \& be \sout{faitful} faithful untill I come even so amen%
					\footnote{%
						\label{EMS-1830-JS-CIR-AUG-1830-NOTE-expansion}%
						This section is greatly expanded upon in the 1835 DC. See notes at \hyperref[DC27]{DC 27}.
						}
				
		% -----
		\subsubsection{Letter to Newel Knight and the Church in Colesville, 28 August 1830}
		% -----
			\label{EMS-1830-JS-28-AUG-1830}
			% \label{EMS-1830-JS-DAY-3LETTERMONTH-YEAR}
		
			\begin{description}
				\item[Print Sources]
			\end{description}

				\begin{tabular}{p{0.5in}p{1in}p{0.5in}p{1in}}
				JSP	 	& D1:172--177	& JSR 		& --- \\
				DC	 	& ---			& HC 		& --- \\
				HDDC 	& ---			& TCHC 		& --- \\
				\end{tabular}
				% HDDC = woodford's DC diss; JSR = Marquardt's JS Revelations; TCHC = Vogel HC
			
			\begin{description}
				\item[Online Access] \href{https://www.josephsmithpapers.org/paper-summary/letter-to-newel-knight-and-the-church-in-colesville-28-august-1830/1}{Here}
				% \item[Analysis] \hyperref[XX]{Here}
			\end{description}
			
			\begin{description}
				\item[Background]
			\end{description}
			
				% It might also be useful to give a two sentence context of the period: e.g., "This speech was given in Nauvoo to a small congregation one month prior to the destruction of the Nauvoo Expositor. These and other tensions may have been on the prophet's mind when he gave the speech."
			
			\begin{description}
				\item[Original Source]
			\end{description}
			
				% brief description of original source material, with reference to JSP or somewhere else for more info
				% Background material: it might be helpful to have a note that says whether a given document was presented as a speech, was written in a journal, etc.
				% Also a comment here about how reliable the source might be, or what shape it is in, if there are large omissions or parts that are hard to understand, and so on.
			
			\begin{description}
				\item[Text]
			\end{description}
			
					\hypertarget{EMS-1830-JS-28-AUG-1830-a}%
				Dearly
					\marginnote{a}
				beloved in the Lord
				
				We are under necessity to disappoint you this time for reasons which I shall mention hereafter, but trusting that your meeting may not be an unprofitable one May you all realize the necessity of getting together often to pray and supplicate at the Throne of grace, that the spirit of the Lord may always rest upon you. Remember that without asking we can receive nothing, therefore ask in faith, and ye shall receive such blessings as God sees fit to bestow upon you.
					\hypertarget{EMS-1830-JS-28-AUG-1830-b}%
				Pray
					\marginnote{b}%
				not with covetous hearts that ye may consume it upon your lusts, but pray earnestly for the best gifts--- fight the good fight of faith, that ye may gain the Crown which is laid up \sout{in} for those that endure faithful unto the end of their probation. Therefore hold fast that which ye have received so liberally from the hand of God so that when the time of refreshing shall come ye may not have labored in vain, but that ye may rest from all your labors and have fulness of joy in the Kingdom of God.
				
					\hypertarget{EMS-1830-JS-28-AUG-1830-c}%
				...Were 
					\marginnote{b}
				it not for the prayers of you few, the Almighty would have thundered down his wrath upon the inhabitants of that place; but be not faint, the day of your deliverance is not far distant, for the judgements of the Lord are already abroad in the earth, and the cold hand of death, will soon pass through your neighborhood, and sweep away some of your most bitter enemies, for you need not suppose that God will be mocked at, and his commandments be trampled under their feet in such a manner as your enemies do, without visiting them in his wrath when they are fully ripe, and behold the angel cries, thrust in your sickle for the harvest is fully ripe; and the earth will soon be reaped.---
					\hypertarget{EMS-1830-JS-28-AUG-1830-d}%
				that
					\marginnote{d}%
				is, the wicked must soon be destroyed from off the face of the earth, for the Lord hath spoken it, and who can stay the hand of the Lord, or who is there that can measure arms with the Almighty, for at his commands the heavens and the earth must pass away, for the day is fast hastening on when the restoration of all things shall be fulfilled, which all the Holy Prophets have prophecid of even unto the gathering in of the House of Israel. Then shall come to pass that the lion shall lie down with the lamb \&c.
					\hypertarget{EMS-1830-JS-28-AUG-1830-e}%
				But
					\marginnote{e}%
				brethren be not dis-couraged when we tell you of perilous times, for they must shortly come, for the sword, famine, and pestilence <are> approaching, for there shall be great destructions upon the face of this land, for ye need not suppose that one jot or tittle of the prophecies of all the Holy Prophets shall fail, and there are many that remain to be fulfilled yet, and the Lord hath said that a short work will he make of it, and the righteous shall be saved if it be as by fire.
				
				``May the grace of God the Father, Son and Holy Ghost be and abide with you from henceforth and forever, Amen.
				
		% -----
		\subsubsection{Revelation, September 1830--A [D\&C 29]}
		% -----
			\label{EMS-1830-JS-SEP-1830-A}
			% \label{EMS-1830-JS-DAY-3LETTERMONTH-YEAR}
		
			\begin{description}
				\item[Print Sources]
			\end{description}

				\begin{tabular}{p{0.5in}p{1in}p{0.5in}p{1in}}
				JSP	 	& D1:177--182	& JSR 		& 80--83	 \\
				DC	 	& Section 29	& HC 		& 1:111--115 \\
				HDDC 	& 415--431		& TCHC 		& 1:78--80 \\
				\end{tabular}
				% HDDC = woodford's DC diss; JSR = Marquardt's JS Revelations; TCHC = Vogel HC
			
			\begin{description}
				\item[Online Access] \href{https://www.josephsmithpapers.org/paper-summary/revelation-september-1830-a-dc-29/1}{Here}
				% \item[Analysis] \hyperref[DC29]{Here}
			\end{description}
			
			\begin{description}
				\item[Background]
			\end{description}
			
				% It might also be useful to give a two sentence context of the period: e.g., "This speech was given in Nauvoo to a small congregation one month prior to the destruction of the Nauvoo Expositor. These and other tensions may have been on the prophet's mind when he gave the speech."
			
			\begin{description}
				\item[Original Source]
			\end{description}
			
				% brief description of original source material, with reference to JSP or somewhere else for more info
				% Background material: it might be helpful to have a note that says whether a given document was presented as a speech, was written in a journal, etc.
				% Also a comment here about how reliable the source might be, or what shape it is in, if there are large omissions or parts that are hard to understand, and so on.
			
			\begin{description}
				\item[Text]
			\end{description}
			
				\begin{tightcenter}
				29\supers{th} Commandment AD 1830
				\end{tightcenter}
				
				A Revelation to Six Elders of the Church \& three members they understood from Holy Writ that the time had come <that> the People of God should see eye to eye \& they seeing somewhat different upon the death of Adam (that is his transgression) therefor they made it a subject of Prayer \& enquired of the Lord \& thus came the word of the Lord through Joseph the seer <saying given> At Fayette Seneca County State of New York
				
				Listen to the voice of Jesus christ your Redeemer the great I am whose arm of mercy hath atoned for your sins who will gether his People even as a hen gethereth her Chickens under her wings even as many as will hearken to my voice \& humble themselves before me \& call upon me in mighty prayer Behold Verily Verily I say unto you at this time your sins are forgiven you Therefore ye Receive these things but remember to sin no more lest perils shall come upon you Verily I say unto you that ye are chosen out of the World to declare my Gospel with the sound of Rejoiceing as with the voice of a Trump lift up your hearts \& be glad for I am in your midst \& am your advocate with the Father \& it is his good will to give you the kingdom \& as it is written Watsoever ye shall ask in faith being united in prayer according to my command ye shall receive \& ye are called to bring to pass the gethering of mine Elect for mine Elect hear my voice \& harden not their hearts wherefore the decree hath gone forth from the father that they shall be gethered in unto one place upon the the face of this land to prepare their Hearts \& be prepared in all things against the day of tribulation \& desolation is sent forth upon the wicked for the hour is nigh \& the day soon at hand when the Earth is ripe \& all the proud \& they that do wickedly shall be as stuble \& I will burn them up that wickedness shall not be upon the Earth for the hour is nigh \& the \sout{day soon at hand} which was spoken by mine Apostles must be fulfilled for as they spoke so shall it come to pass for I will reveal myself from Heaven with Power \& great glory with all the hosts thereof \& dwell in righteousness with men on Earth a thousand Years \& the wicked shall not stand \& again Verily Verily I say unto you \& it hath gone forth in a firm decree by the will of the father that mine Apostles the twelve which were with me in my ministery at Jerusalem shall stand at my right hand at the day of my comeing in a piller of fire being clothed with robes of righteousness with crowns upon their heads in glory even as I am to Judge the whole House of Israel even as many as have loved me \& kept my commandments \& none else for a Trump Shall sound both long \& loud even as upon mount Sinia [Sinai] \& all the Earth shall quake \& they shall come forth yea even the dead which died in me to receive a Crown of righteousness \& to be Clothed upon even as I am to be with me that we may be one, but, Behold I say unto you that before this great day shall come the Sun shall be darkened \& the moon shall be turned into blood \& some stars shall fall from Heaven \& there shall be greater signs in the Heaven above \& in the Earth beneath \& there shall be weeping \& waileing among the host of men \& there shall be be a great hailstorm sent forth to destroy the Crops of the Earth \& it shall come to pass because of the wickedness of the World that I will take vengeance upon the Wicked for they will not Repent for the cup of mine indignation is full for Behold my blood shall not cleanse them if they repent not Wherefore I will send <forth> flies upon the face of the Earth which shall take hold of the inhabitants thereof \& shall eat their flesh \& shall cause magots to come in upon them \& their tongues shall be stayed that they shall not utter against me \& their flesh shall fall from off their Bones \& their eyes from their sockets \& it shall come to pass that \sout{their} the Beasts of the forest \& the fowls of the air shall devour them up \& that great \& abominable Church which is the whore of all the Earth shall be cast down by devouring fire according as it was spoken by the mouth of Ezekiel the Prophet which spoke of these things which have not come to pass as yet but shurely must as I live for abominations shall not reign \& again Verily Verily I say unto you that when the thousand years are ended \& men again begin to deny their god then will I spare the Earth but for a little Season \& then the end shall come \& the Heaven \& the Earth shall be consumed \& pass away \& there shall be a New Heaven \& a New Earth for all old things shall pass away \& all things shall become New even the Heaven \& the Earth \& all the fulness thereof both men \& beasts the fowls of the air \& the fishes of the Sea \& not one hair neither moat [mote] shall be lost for it is the workmanship of mine hand But Verily I say unto you before the Earth shall pass away Michael mine Archangel shall sound his trump \& then shall all the dead awake for their graves shall be opened \& they shall come forth yea even all \& the righteous shall be gethered on my right hand unto eternal life \& the wicked on my left hand will I be ashamed to own before the father Wherefore I will say unto them depart from me ye cursed into everlasting fire prepared for the devil \& his Angels \& now Behold I say unto you never at any time have I declared from mine own mouth that they should return for where I am they cannot come for they have no power but remember that all my Judgements are not given unto men \& as the words have gone forth out of my mouth even so shall they be fulfilled that the first shall be last \& that the last shall be first in all things Whatsoever I have created by the word of my Power which is the Power of my spirit for by the Power of my Spirit created I them yea all things both Spiritual \& Temperally firstly spiritual secondly temporally which is the Begining of my work \& again firstly temporal \& secondly spiritual which is the last of my work speaking unto you that ye may naturally understand but unto myself my work hath no end neither begining But it is given unto you that ye may understand because ye have asked it of me \& are agreed Wherefore Verily I say unto you that all things unto me are Spiritual \& not at any time have I given unto you a law which was temporal neither any man nor the childern of men Neither Adam your father whom I created Behold I \sout{give} <gave> unto him that he should be an agent unto himself \& I gave unto him a commandment but no temporal Commandment gave I unto him for my commandments are spiritual they are not Natural nor temporal neither carnal nor sensual \& it came to pass that Adam being tempted of the Devil for Behold the Devil was before Adam for he rebelled against me saying give me thine honour which is my Power \& also a third part of the host of Heaven turned he away from me Because of their agency \& they were thrust down \& thus came the Devil \& his Angels \& Behold a place prepared for them which place is Hell \& it \sout{came to pass} Must needs be that the Devil should tempt the children of men or they could not be agents unto themselves for if they never should have bitter they could not k[n]ow the Sweet Wherefore it came to pass that the Devil tempted Adam \& he partook of the forbiden fruit \& transgressed the commandment wherein he became subject to the will of the Devil Because he yielded unto temptation Wherefore I the Lord God caused that he should be cast out from the Garden of Edan from my presence because of his transgression Wherein he became spiritually dead which \sout{death} is the first death even that same death which is the last death which is spiritual \sout{which shall be} pronounced upon the wicked which shall be when I shall say depart ye Cursed But Behold I say unto you that I the Lord God gave unto Adam \& unto his seed that they should not Die as to the temporal death untill I the Lord God should send forth Angels to declare unto them Repentance \& redemption through faith on the name of mine only begotten Son \& thus did I the Lord God appoint unto man the days of his probation that by his natural death he might be raised in immortality unto eternal life even as many as would believe on my name \& they that believe not unto eternal damnation for they cannot be redeemed from their spiritual fall Because they repent not for they love darkness more than light \& their deeds are evil \& they receive their wages of whom they list to obey But Behold I say unto you that little children are redeemed from the foundation of the \sout{word} world through mine only begotten Wherefore they cannot sin for power is not given to Satan to tempt little children until they begin to be accountable before me for it is given unto them even as I will according to mine own \sout{will} pleasure that great things may be required at the hand of their fathers \& again I say unto you that whoso having knowledge have not I commanded to Repent \& he that hath no understanding it remaineth in me to do according as it is written \& now behold I declare no more unto you at this time amen
				
		% -----
		\subsubsection{Revelation, September 1830--B [D\&C 28]}
		% -----
			\label{EMS-1830-JS-SEP-1830-B}
			% \label{EMS-1830-JS-DAY-3LETTERMONTH-YEAR}
		
			\begin{description}
				\item[Print Sources]
			\end{description}

				\begin{tabular}{p{0.5in}p{1in}p{0.5in}p{1in}}
				JSP	 	& D1:183--186	& JSR 		& 83--89\footnote{Also a lot of commentary here.} \\
				DC	 	& Section 28	& HC 		& 1:110--111 \\
				HDDC 	& 404--414		& TCHC 		& 1:78 \\
				\end{tabular}
				% HDDC = woodford's DC diss; JSR = Marquardt's JS Revelations; TCHC = Vogel HC
			
			\begin{description}
				\item[Online Access] \href{https://www.josephsmithpapers.org/paper-summary/revelation-september-1830-b-dc-28/1}{Here}
				% \item[Analysis] \hyperref[DC28]{Here}
			\end{description}
			
			\begin{description}
				\item[Background]
			\end{description}
			
				% It might also be useful to give a two sentence context of the period: e.g., "This speech was given in Nauvoo to a small congregation one month prior to the destruction of the Nauvoo Expositor. These and other tensions may have been on the prophet's mind when he gave the speech."
			
			\begin{description}
				\item[Original Source]
			\end{description}
			
				% brief description of original source material, with reference to JSP or somewhere else for more info
				% Background material: it might be helpful to have a note that says whether a given document was presented as a speech, was written in a journal, etc.
				% Also a comment here about how reliable the source might be, or what shape it is in, if there are large omissions or parts that are hard to understand, and so on.
			
			\begin{description}
				\item[Text]
			\end{description}
			
				\begin{tightcenter}
				30 Commandment AD 1831 [1830]
				\end{tightcenter}
				
				A Revelation to Oliver [Cowdery] his Call to the Lamanitse \&c given at Fayette Seneca County State of New York
				
				Behold I say unto you Oliver that it shall be given thee that thou shalt be heard by the Church in all things Whatsoever thou shalt teach <them> by the Comforter concerning the Revelations \& commandments which I have given But Behold Verily Verily I say unto you no one shall be appointed to Receive commandments \& Revelations in this Church excepting my Servent Joseph for he Receiveth them even as Moses \& thou shalt be obedient unto the things which I shall give unto him Even as Aaaron to declare faithfully the commandments \& the Revelations with power \& authority unto the Church \& if thou art led at any time by the comforter to speak or teach or at all times by the way of Commandment unto the Church thou mayest do it But thou shalt not write by way of Commandment \sout{unto the Church} but by wisdom \& thou shalt not command him which is at thy head \& at the head of the Church for I have given him the keys of the mysteries of the Revelations which are sealed until I shall appoint unto him another in his stead \& now Behold I say unto you that thou shalt go unto the Lamanites \& Preach my Gospel unto them \& cause my Church to be established among them \& thou shalt have Revelations but write them not by the way of Commandment \& Now Behold I say unto you that it is not Revealed \& no man knoweth where the City shall be built But it shall be given hereafter Behold I say unto you that it shall be among the Lamanites[.] thou shalt not l[e]ave this place until after the Conference \& my servent Joseph shall be appointed to rule the conference by the voice of it \& what he saith to thee that thou shalt tell And again thou shalt take thy Brother Hyram [Hiram Page] Between him \& thee alone \& tell him that those things which he hath written from that Stone are not of me \& that Satan deceiveth him for Behold those things have not been appointed unto him Neither shall any thing be appointed unto any of this Church contrary to the Church Articles \& Covenants for all things must be done in order \& by Common consent in the Church by the prayer of faith \& thou shalt settle all these things according to the Covenants of the Church before thou shalt take thy Journey among the Lamanites \& it shall be given thee from the time that thou shalt go until the time that thou shalt return what thou shalt do \& thou must open thy mouth at all times declaring my Gospel with the sound of Rejoiceing even so amen
				
		% -----
		\subsubsection{Revelation, September 1830--C [D\&C 30:1--4]}
		% -----
			\label{EMS-1830-JS-SEP-1830-C}
			% \label{EMS-1830-JS-DAY-3LETTERMONTH-YEAR}
		
			\begin{description}
				\item[Print Sources]
			\end{description}

				\begin{tabular}{p{0.5in}p{1in}p{0.5in}p{1in}}
				JSP	 	& D1:187--188	& JSR 		& 89	 \\
				DC	 	& Section 30	& HC 		& 1:116 \\
				HDDC 	& 432--444		& TCHC 		& 1:81 \\
				\end{tabular}
				% HDDC = woodford's DC diss; JSR = Marquardt's JS Revelations; TCHC = Vogel HC
			
			\begin{description}
				\item[Online Access] \href{https://www.josephsmithpapers.org/paper-summary/revelation-september-1830-c-dc-301-4/1}{Here}
				% \item[Analysis] \hyperref[DC30]{Here}
			\end{description}
			
			\begin{description}
				\item[Background]
			\end{description}
			
				% It might also be useful to give a two sentence context of the period: e.g., "This speech was given in Nauvoo to a small congregation one month prior to the destruction of the Nauvoo Expositor. These and other tensions may have been on the prophet's mind when he gave the speech."
			
			\begin{description}
				\item[Original Source]
			\end{description}
			
				% brief description of original source material, with reference to JSP or somewhere else for more info
				% Background material: it might be helpful to have a note that says whether a given document was presented as a speech, was written in a journal, etc.
				% Also a comment here about how reliable the source might be, or what shape it is in, if there are large omissions or parts that are hard to understand, and so on.
			
			\begin{description}
				\item[Text]
			\end{description}
			
				\begin{tightcenter}
				31 Commandment AD 1830
				\end{tightcenter}
				
				A commandment to David [Whitmer] tellilg [telling] him that he feared man more than god \&c given at fayette Seneca County New York
				
				Behold I say unto you David that thou hast feared man \& hast not relyed upon me for strength as thou hast ought But thy mind has been on the things of Earth more than on the things of me thy Maker \& the ministery whereunto thou hast been called \& thou hast not given heed unto my Spirit \& to those who were set over thee But hast been persuaded by those whom I have not commanded wherefore thou art left to enquire for thy self at my hand \& ponder upon the things which you have Received \& thy home shall be at thy fathers house until I give unto thee further commandment \& thou shalt attend to the ministery in the Church \& before the world \& in these regions round about amen
				
		% -----
		\subsubsection{Revelation, September 1830--D [D\&C 30:5--8]}
		% -----
			\label{EMS-1830-JS-SEP-1830-D}
			% \label{EMS-1830-JS-DAY-3LETTERMONTH-YEAR}
		
			\begin{description}
				\item[Print Sources]
			\end{description}

				\begin{tabular}{p{0.5in}p{1in}p{0.5in}p{1in}}
				JSP	 	& D1:188--189	& JSR 		& 89--90 \\
				DC	 	& Section 30	& HC 		& 1:116 \\
				HDDC 	& 432--444		& TCHC 		& 1:81 \\
				\end{tabular}
				% HDDC = woodford's DC diss; JSR = Marquardt's JS Revelations; TCHC = Vogel HC
			
			\begin{description}
				\item[Online Access] \href{https://www.josephsmithpapers.org/paper-summary/revelation-september-1830-d-dc-305-8/1}{Here}
				% \item[Analysis] \hyperref[DC30]{Here}
			\end{description}
			
			\begin{description}
				\item[Background]
			\end{description}
			
				% It might also be useful to give a two sentence context of the period: e.g., "This speech was given in Nauvoo to a small congregation one month prior to the destruction of the Nauvoo Expositor. These and other tensions may have been on the prophet's mind when he gave the speech."
			
			\begin{description}
				\item[Original Source]
			\end{description}
			
				% brief description of original source material, with reference to JSP or somewhere else for more info
				% Background material: it might be helpful to have a note that says whether a given document was presented as a speech, was written in a journal, etc.
				% Also a comment here about how reliable the source might be, or what shape it is in, if there are large omissions or parts that are hard to understand, and so on.
			
			\begin{description}
				\item[Text]
			\end{description}
			
				\begin{tightcenter}
				32\supers{nd} Commandment AD 1830
				\end{tightcenter}
				
				A Revelation to Peter [Whitmer Jr.] his calling to the Lamanites \&c given at Fayette Seneca County state of New York
				
				Behold I say unto you Peter that thou shalt take thy Journey with thy Brother oliver [Cowdery] for the time has come that it is expedient in me that thou shalt open thy mouth to declare my Gospel Therefore fear not but give heed unto the words \& advice of thy Brother which he shall give thee \& be thou afflicted in all his afflictions ever lifting thy heart up unto me in prayer \& faith for thine \& his deliverance for I have given unto him to build my Church among thy Brethren the Lamanites \& none have I appointed to be over him in the Church except it is his Brother Joseph wherefore give heed unto \sout{those} <these> things \& be dilligent in keeping my commandments \& thou shalt be blessed unto eternal life \& thus it is amen---
				
		% -----
		\subsubsection{Revelation, September 1830--E [D\&C 30:9--11]}
		% -----
			\label{EMS-1830-JS-SEP-1830-E}
			% \label{EMS-1830-JS-DAY-3LETTERMONTH-YEAR}
		
			\begin{description}
				\item[Print Sources]
			\end{description}

				\begin{tabular}{p{0.5in}p{1in}p{0.5in}p{1in}}
				JSP	 	& D1:189--190	& JSR 		& 90	 \\
				DC	 	& Section 30	& HC 		& 1:116 \\
				HDDC 	& 432--444		& TCHC 		& 1:81 \\
				\end{tabular}
				% HDDC = woodford's DC diss; JSR = Marquardt's JS Revelations; TCHC = Vogel HC
			
			\begin{description}
				\item[Online Access] \href{https://www.josephsmithpapers.org/paper-summary/revelation-september-1830-e-dc-309-11/1}{Here}
				% \item[Analysis] \hyperref[DC30]{Here}
			\end{description}
			
			\begin{description}
				\item[Background]
			\end{description}
			
				% It might also be useful to give a two sentence context of the period: e.g., "This speech was given in Nauvoo to a small congregation one month prior to the destruction of the Nauvoo Expositor. These and other tensions may have been on the prophet's mind when he gave the speech."
			
			\begin{description}
				\item[Original Source]
			\end{description}
			
				% brief description of original source material, with reference to JSP or somewhere else for more info
				% Background material: it might be helpful to have a note that says whether a given document was presented as a speech, was written in a journal, etc.
				% Also a comment here about how reliable the source might be, or what shape it is in, if there are large omissions or parts that are hard to understand, and so on.
			
			\begin{description}
				\item[Text]
			\end{description}
			
				\begin{tightcenter}
				33 Commandment AD 183[0]
				\end{tightcenter}
				
				A Revelation to John [Whitmer] his call to the Ministery \sout{give} \&c given at Fayette Seneca County State of New York
				
				Behold I say unto John that thou shalt commence from this time forth to proclaim my Gospel as with the voice of a Trump \& thy Labour shall be at thy Brother Philip\supersu{s}\supers{.} [Philip Burroughs's] \& in that region round about yea wheresoever thou canst be heard until I command thee to go from hence \& thy whole Labour shall be in my Zion with all thy Soul from henceforth yea thou shalt ever open thy mouth in my cause not fearing what man can do for I a$\diamond$ with thee even so amen
				
		% -----
		\subsubsection{Revelation, September 1830--F [D\&C 31]}
		% -----
			\label{EMS-1830-JS-SEP-1830-F}
			% \label{EMS-1830-JS-DAY-3LETTERMONTH-YEAR}
		
			\begin{description}
				\item[Print Sources]
			\end{description}

				\begin{tabular}{p{0.5in}p{1in}p{0.5in}p{1in}}
				JSP	 	& D1:193--195	& JSR 		& 91	 \\
				DC	 	& Section 31	& HC 		& 1:116--117 \\
				HDDC 	& 432--444		& TCHC 		& 1:81--82 \\
				\end{tabular}
				% HDDC = woodford's DC diss; JSR = Marquardt's JS Revelations; TCHC = Vogel HC
			
			\begin{description}
				\item[Online Access] \href{https://www.josephsmithpapers.org/paper-summary/revelation-september-1830-f-dc-31/1}{Here}
				% \item[Analysis] \hyperref[DC31]{Here}
			\end{description}
			
			\begin{description}
				\item[Background]
			\end{description}
			
				% It might also be useful to give a two sentence context of the period: e.g., "This speech was given in Nauvoo to a small congregation one month prior to the destruction of the Nauvoo Expositor. These and other tensions may have been on the prophet's mind when he gave the speech."
			
			\begin{description}
				\item[Original Source]
			\end{description}
			
				% brief description of original source material, with reference to JSP or somewhere else for more info
				% Background material: it might be helpful to have a note that says whether a given document was presented as a speech, was written in a journal, etc.
				% Also a comment here about how reliable the source might be, or what shape it is in, if there are large omissions or parts that are hard to understand, and so on.
			
			\begin{description}
				\item[Text]
			\end{description}
			
				\begin{tightcenter}
				34\supers{th} Commandment AD 1830
				\end{tightcenter}
				
				A Revelation to Thomas [B. Marsh] his call to the ministiry \&c gaven at Fayette Seneca County State of New York
				
				Thomas my Son Blessed art thou Be[c]ause of thy faith in my words Behold thou hast had many afflictions because of thy family Nevertheless I will bless thee \& thy family yea thy little ones \& the day cometh that they will believe \& know the truth \& be one with thee in my Church lift up your heart \& rejoice for the hour of your \sout{Mishion} mission is come \& thy tongue shall be loosed \& thou shalt declare glad tidings of great joy unto this generation thou shalt declare the things which have been revealed unto my Servent Joseph thou shalt begin to preach from this time forth yea to Reap in the field which is white already to be burned Therefore thrust in thy Sickle with all thy Soul \& thy sins are forgiven thee \& thou shalt be laden with sheaves upon thy Back for the labourer is worthy of his hire Wherefore thy family shall live Behold Verily I say unto you go from them only for a little time \& declare my word \& I will prepare a place for them yea I will open the hearts of the People \& they will Receive thee \& I will establish a church by thy hand \& thou shalt strengthen them \& prepare them against the time when the gethering shall be[.] be patient in afflictions \& sufferings revile not against Those that revile govern thy house in meekness \& be steadfast Behold I say unto you that thou shalt be a P[h]ysician unto the Church but not unto the World for they will not receive thee go thy way whithersoever I will \& it shall be given thee by the Comforter what thou shalt do \& whither thou shalt go pray always lest ye enter into temptation \& loose thy reward be faithful unto the end \& Lo! I am with you these words are not of man neither of men but of me even Jesus Christ your Redeemer by the will of the father even so amen
				
		% -----
		\subsubsection{Revelation, October 1830--A [D\&C 32]}
		% -----
			\label{EMS-1830-JS-OCT-1830-A}
			% \label{EMS-1830-JS-DAY-3LETTERMONTH-YEAR}
		
			\begin{description}
				\item[Print Sources]
			\end{description}

				\begin{tabular}{p{0.5in}p{1in}p{0.5in}p{1in}}
				JSP	 	& D1:200--202	& JSR 		& 92	 \\
				DC	 	& Section 32	& HC 		& 1:118--120 \\
				HDDC 	& 445--452		& TCHC 		& 1:83 \\
				\end{tabular}
				% HDDC = woodford's DC diss; JSR = Marquardt's JS Revelations; TCHC = Vogel HC
			
			\begin{description}
				\item[Online Access] \href{https://www.josephsmithpapers.org/paper-summary/revelation-october-1830-a-dc-32/1}{Here}
				% \item[Analysis] \hyperref[DC32]{Here}
			\end{description}
			
			\begin{description}
				\item[Background]
			\end{description}
			
				% It might also be useful to give a two sentence context of the period: e.g., "This speech was given in Nauvoo to a small congregation one month prior to the destruction of the Nauvoo Expositor. These and other tensions may have been on the prophet's mind when he gave the speech."
			
			\begin{description}
				\item[Original Source]
			\end{description}
			
				% brief description of original source material, with reference to JSP or somewhere else for more info
				% Background material: it might be helpful to have a note that says whether a given document was presented as a speech, was written in a journal, etc.
				% Also a comment here about how reliable the source might be, or what shape it is in, if there are large omissions or parts that are hard to understand, and so on.
			
			\begin{description}
				\item[Text]
			\end{description}
			
				Revelation to Parley [P.] Pratt to go to th[e] wilderness
				
				And now concerning my servant Parley behold I say unto him that as I live I will that he shall declare my gospel and Learn of me and be meek and lowly of heart and that which I have appointed unto him is that he shall go with my servant Oliver [Cowdery] and Peter [Whitmer Jr.] into the wilderness among the Lamanites and Ziba [Peterson] also shall go with them and I myself will go with them and be in their midst and I am their advocate with the Father and nothing shall prevail and they shall give heed to that which is writen and pretend to no other revelation and they shall pray always that I may unfold them to their understanding and they shall give heed unto these words and trifle not and I will bless them amen Manchester Oct 1830---
				
		% -----
		\subsubsection{Revelation, October 1830--B [D\&C 33]}
		% -----
			\label{EMS-1830-JS-OCT-1830-B}
			% \label{EMS-1830-JS-DAY-3LETTERMONTH-YEAR}
		
			\begin{description}
				\item[Print Sources]
			\end{description}

				\begin{tabular}{p{0.5in}p{1in}p{0.5in}p{1in}}
				JSP	 	& D1:205--208	& JSR 		& 92--94	 \\
				DC	 	& Section 33	& HC 		& 1:126--127 \\
				HDDC 	& 453--460		& TCHC 		& 1:97--98 \\
				\end{tabular}
				% HDDC = woodford's DC diss; JSR = Marquardt's JS Revelations; TCHC = Vogel HC
			
			\begin{description}
				\item[Online Access] \href{https://www.josephsmithpapers.org/paper-summary/revelation-october-1830-b-dc-33/1}{Here}
				% \item[Analysis] \hyperref[DC33]{Here}
			\end{description}
			
			\begin{description}
				\item[Background]
			\end{description}
			
				% It might also be useful to give a two sentence context of the period: e.g., "This speech was given in Nauvoo to a small congregation one month prior to the destruction of the Nauvoo Expositor. These and other tensions may have been on the prophet's mind when he gave the speech."
			
			\begin{description}
				\item[Original Source]
			\end{description}
			
				% brief description of original source material, with reference to JSP or somewhere else for more info
				% Background material: it might be helpful to have a note that says whether a given document was presented as a speech, was written in a journal, etc.
				% Also a comment here about how reliable the source might be, or what shape it is in, if there are large omissions or parts that are hard to understand, and so on.
			
			\begin{description}
				\item[Text]
			\end{description}
			
				\begin{tightcenter}
				35 \& 36\supers{th} \sout{Commands} Commandment AD 1830
				\end{tightcenter}
				
				A Commandment to Ezra [Thayer] \& Northrop [Sweet] th[e]ir call to the ministery \&c given at Fayette Seneca County State of New York
				
				Saying I say unto you Ezra \& Northrop open ye your ears \& hearken to the voice of the Lord your God whose word is quick \& powerfull sharper than a twoedged sword to the dividing asunder of the Joints \& marrow Soul \& spirit \& is a decerner of the thoughts \& intents of the heart for Verily Verily I say unto you that ye are called to lift up voices as with the sound of a Trump to declare my Gospel unto a Crooked \& a perverse generation for Behold the field is white already to harvest \& it is the Elvenenth hour \& for the last time that I shall call labouerers into my vineyard \& my vineyard has become corrupted evry whit \& there is none that doeth good save it is a few only <\& they do err in many instances because of Priest crafts> all having corrupt minds \& Verily Verily I say unto you that this Church have I established \& called forth out of the Wilderness \& even so will I gether mine \sout{elect} elect from the four quarters of the Earth even as many as will believe \sout{in} <on> my name \& hearken unto my voice yea Verily Verily I say unto you that the field is white already to harvest Wherefore thrust in thy sickles \& reap with all thy might\sout{s} mind \& strength open thy mouth \& it shall be filled \& thou shalt become even as Nephi of old who Journ[ey]ed from Jerusalem in the wilderness yea open thy mouth \& spare not \& thou shalt be laden with Sheaves upon thy Back for lo I am with \sout{thee} you yea Open thy mouth \& it shall be filled saying Repent repent \& prepare ye the way of the Lord \& make his path\sout{s} strait for the Kingdom of Heaven is at hand yea Repent \& be Baptized evry one of you for the remission of sins yea be baptized even by water \& then cometh the Baptism of fire \& the Holy Ghost Behold Verily Verily I say unto you this is my Gospel \& Remember they shall have faith in me or they can in no wise be saved \& upon this Rock I will build my Church yea upon this Rock ye are built \& the gaits of Hell shall not prevail against you \& ye shall remember the Church Articles \& Covenants to keep them \& whoso haveing faith ye shall confirm in my Church by the laying on of the hands \& I will bestow the gift of the Holy Ghost upon them \& the Book of Mormon \& the Bible \sout{is} are given of me for thine instructio[n] \& the power of my spirit quickeneth all things Wherefore be faithful praying always having your lamps trimmed \& burning \& oil with you that ye may be Ready at the coming of the Bride groom for Behold Verily Verily I say unto you that I come quickly even so amen
				
		% -----
		\subsubsection{Revelation, 4 November 1830 [D\&C 34]}
		% -----
			\label{EMS-1830-JS-4-NOV-1830}
			% \label{EMS-1830-JS-DAY-3LETTERMONTH-YEAR}
		
			\begin{description}
				\item[Print Sources]
			\end{description}

				\begin{tabular}{p{0.5in}p{1in}p{0.5in}p{1in}}
				JSP	 	& D1:208--211	& JSR 		& 94	 \\
				DC	 	& Section 34	& HC 		& 1:128 \\
				HDDC 	& 461--466		& TCHC 		& 1:98 \\
				\end{tabular}
				% HDDC = woodford's DC diss; JSR = Marquardt's JS Revelations; TCHC = Vogel HC
			
			\begin{description}
				\item[Online Access] \href{https://www.josephsmithpapers.org/paper-summary/revelation-4-november-1830-dc-34/1}{Here}
				% \item[Analysis] \hyperref[DC34]{Here}
			\end{description}
			
			\begin{description}
				\item[Background]
			\end{description}
			
				% It might also be useful to give a two sentence context of the period: e.g., "This speech was given in Nauvoo to a small congregation one month prior to the destruction of the Nauvoo Expositor. These and other tensions may have been on the prophet's mind when he gave the speech."
			
			\begin{description}
				\item[Original Source]
			\end{description}
			
				% brief description of original source material, with reference to JSP or somewhere else for more info
				% Background material: it might be helpful to have a note that says whether a given document was presented as a speech, was written in a journal, etc.
				% Also a comment here about how reliable the source might be, or what shape it is in, if there are large omissions or parts that are hard to understand, and so on.
			
			\begin{description}
				\item[Text]
			\end{description}
			
				\begin{tightcenter}
				37\supers{th} Commandment AD 1830
				\end{tightcenter}
				
				A Commandment to Orson [Pratt] his call to the ministery \&c given at Fayette Seneca county State of New York
				
				My Son Orson hearken ye \& Behold what I the Lord God say unto you even Jesus Christ your Redeemer the light \& the life of the world \sout{\&} a light which shineth in darkness \& the darkness Comprehendeth it not who so loved the world that he gave his own life that as many as would believe might become the Sons \& daughter[s] of God Wherefore ye are my Son \& blessed are ye because ye have believed \& more blessed are ye because ye are called of me to Preach my Gospel to lift up your voice as with \sout{a} the sound of a Trump both long \& loud \& cry repentance to a crooked \& perverse generation prepareing the way of the Lord for his second Coming for Behold Verily Verily I say unto you the time is soon at hand that I will come in a cloud with power \& great glory \& it shall be a great day at the time of my coming for all nations shall tremble but before that great day shall come the sun shall be darkened \& the moon be turned into blood \& the stars shall refuse their shineing \& some shall fall \& great distructions await the wicked Wherefore lift up thy voice \& spare not for the Lord God hath spoken therefore Prophecy \& it shall be given by the power of the Holy Ghost \& if ye are faithful behold I am with you until I come \& Verily Verily I say unto you I come quickly even so your Lord \& your redeemer amen
				
		% -----
		\subsubsection{Letter to the Church in Colesville, 2 December 1830}
		% -----
			\label{EMS-1830-JS-2-DEC-1830}
			% \label{EMS-1830-JS-DAY-3LETTERMONTH-YEAR}
		
			\begin{description}
				\item[Print Sources]
			\end{description}

				\begin{tabular}{p{0.5in}p{1in}p{0.5in}p{1in}}
				JSP	 	& D1:214--219	& JSR 		& --- \\
				DC	 	& ---			& HC 		& --- \\
				HDDC 	& ---			& TCHC 		& --- \\
				\end{tabular}
				% HDDC = woodford's DC diss; JSR = Marquardt's JS Revelations; TCHC = Vogel HC
			
			\begin{description}
				\item[Online Access] \href{https://www.josephsmithpapers.org/paper-summary/letter-to-the-church-in-colesville-2-december-1830/1}{Here}
				% \item[Analysis] \hyperref[XX]{Here}
			\end{description}
			
			\begin{description}
				\item[Background]
			\end{description}
			
				% It might also be useful to give a two sentence context of the period: e.g., "This speech was given in Nauvoo to a small congregation one month prior to the destruction of the Nauvoo Expositor. These and other tensions may have been on the prophet's mind when he gave the speech."
			
			\begin{description}
				\item[Original Source]
			\end{description}
			
				% brief description of original source material, with reference to JSP or somewhere else for more info
				% Background material: it might be helpful to have a note that says whether a given document was presented as a speech, was written in a journal, etc.
				% Also a comment here about how reliable the source might be, or what shape it is in, if there are large omissions or parts that are hard to understand, and so on.
			
			\begin{description}
				\item[Text]
			\end{description}
			
					\hypertarget{EMS-1830-JS-2-DEC-1830-a}%
				Dearly
					\marginnote{a}
				beloved in the Lord
				
				According to your prayers, the Lord hath called, chosen, ordained, sanctified and sent unto you, another servant and Apostle separated unto his gospel through Jesus Christ \sout{his} our Redeemer, to whom be all honor \& praise henceforth and forever--- even our beloved brother Orson Pratt, the bearer of these lines. Whom I recommend unto you as a faithful Servant in the Lord, through Jesus Christ our Redeemer, Amen.
					
					\hypertarget{EMS-1830-JS-2-DEC-1830-b}%
				To
					\marginnote{b}				
				the Church in Colesville---
					
				Having many things to write to you, but being assured that ye are not ignorant of all that I can write to you, finally I would inform you that Zion is prospering here, there are many serious inquirers in this place, who are seeking the Lord. It gave us much joy to hear from you, to hear that God is softening the hearts of the children of men in that place, it being the seat of Satan. But blessed be the name of God, it also hath become the abode of our savior, and may you all be faithful and wait for the time of our Lord, for his appearing is nigh\sout{t} at hand.
					\hypertarget{EMS-1830-JS-2-DEC-1830-c}%
				But
					\marginnote{c}%
				the time, and the season, Brethren, ye have no need that I write unto you, for ye yourselves perfectly know that the day of the Lord so cometh as a thief in the night: for when they shall say peace and safety, then sudden destruction cometh upon them, as travail upon a woman, but they shall not escape. But ye, brethren are not in darkness, therefore let us not sleep as do others, but let us watch and be sober, for they that sleep, sleep in the night, and they that be drunken are drunken in the night, but let us who be of the day, be sober, putting on the breastplate of faith and law, and for a helmet, the hope of salvation. For God hath not appointed us \sout{un}to wrath; but to obtain [sa]lvation through our Lord Jesus Christ.
					\hypertarget{EMS-1830-JS-2-DEC-1830-d}%
				Wherefore
					\marginnote{d}%
				comfort one another, even as ye also do; for perilous times are at hand, for behold the dethronement and deposition of the kings in the eastern continent,--- the whirlwinds in the West India Islands, it has destroyed a number of vessels, uprooted buildings and strewed them in the air; the fields of spices have been destroyed, and the inhabitants have barely escaped with their lives, and many have been buried under the ruins.
					\hypertarget{EMS-1830-JS-2-DEC-1830-e}%
				In
					\marginnote{e}%
				Columbia, South America, they are at war and peace is taken from the earth in part. and it will soon be in whole, yea destructions are at our doors, and they soon will be in the houses of the wicked, and they that know not God. Yea lift up your heads and rejoice for your redemption draweth nigh. We are the most favored people that ever have been from the foundation of the world, if we remain faithful in keeping the commandments of our God.
					\hypertarget{EMS-1830-JS-2-DEC-1830-f}%
				Yea,
					\marginnote{f}%
				even Enoch, the seventh from Adam beheld our day and rejoiced, and the prophets from that day forth have prophecied of the second coming of our Lord and Savior Jesus Christ, and rejoiced at the day of rest of the Saints, Yea, and the Apostle of our Savior also did rejoice in his appear\supersu{$\diamond$} in a cloud with the host of Heaven, to dwell with man on the earth a thousand years. Therefore we have reason to rejoice. Behold the prophecies of the Book of Mormon are fulfilling as fast as time can bring it about. The Spirit of the Living God is upon me therefore who will say that I shall not prophecy.
					\hypertarget{EMS-1830-JS-2-DEC-1830-g}%
				The 
					\marginnote{g}%
				time is soon at hand that we shall have to flee whithersoever the Lord will, for safety, Fear not those who are making you an offender for a word but be faithful in witnessing unto a crooked and a perverse generation, that \sout{thy} the day of the coming of our Lord and Savior is at hand. Yea, prepare ye the way of the Lord, make strait his path. Who will shrink because of offences, for offences must come, but woe to them by whom they come, for the rock must fall on them and grind them to powder, for the fulness of the gentiles is come in, and woe will be unto them if they do not repent and be baptized in the name of our Lord and Savior Jesus Christ for the remission of their sins, and come in at the strait gate and be numbered with the House of Israel,
					\hypertarget{EMS-1830-JS-2-DEC-1830-h}%
				for
					\marginnote{h}%
				God will not always be mocked, and not pour out his wrath upon those that blaspheme his holy name, for the sword, famines and destruction will soon overtake them in their wild career, for God will avenge, and pour out his phials of wrath, and save his elect. And all those who will obey his commandments are his elect, and he will soon gather them from the four winds of heaven, from one quarter of the earth to the other, to a place whithersoever he will, therefore in your patience possess ye your souls. Amen.
				
		% -----
		\subsubsection{Revelation, 7 December 1830 [D\&C 35]}
		% -----
			\label{EMS-1830-JS-7-DEC-1830}
			% \label{EMS-1830-JS-DAY-3LETTERMONTH-YEAR}
		
			\begin{description}
				\item[Print Sources]
			\end{description}

				\begin{tabular}{p{0.5in}p{1in}p{0.5in}p{1in}}
				JSP	 	& D1:219--223	& JSR 		& 95--96	 \\
				DC	 	& Section 35	& HC 		& 1:129--131 \\
				HDDC 	& 467--483		& TCHC 		& 1:99--100 \\
				\end{tabular}
				% HDDC = woodford's DC diss; JSR = Marquardt's JS Revelations; TCHC = Vogel HC
			
			\begin{description}
				\item[Online Access] \href{https://www.josephsmithpapers.org/paper-summary/revelation-7-december-1830-dc-35/1}{Here}
				% \item[Analysis] \hyperref[DC35]{Here}
			\end{description}
			
			\begin{description}
				\item[Background]
			\end{description}
			
				% It might also be useful to give a two sentence context of the period: e.g., "This speech was given in Nauvoo to a small congregation one month prior to the destruction of the Nauvoo Expositor. These and other tensions may have been on the prophet's mind when he gave the speech."
			
			\begin{description}
				\item[Original Source]
			\end{description}
			
				% brief description of original source material, with reference to JSP or somewhere else for more info
				% Background material: it might be helpful to have a note that says whether a given document was presented as a speech, was written in a journal, etc.
				% Also a comment here about how reliable the source might be, or what shape it is in, if there are large omissions or parts that are hard to understand, and so on.
			
			\begin{description}
				\item[Text]
			\end{description}
			
				\begin{tightcenter}
				38\supers{th} Commandment AD 1830. Dec.\supersu{m}\supers{.} 7\supers{th}
				\end{tightcenter}
				
				A Commandment to Joseph \& Sidney [Rigdon]. Sidneys call to writing for Joseph \&c
				
				Saying Listen to the voice of the Lord your God even Alpha \& Omega the begining \& the end whose course is one eternal round the same to day as yesterday \& for ever I am Jesus Christ the son of God who was crusified for the sins of the World even as will believe on my name that they may become the sons of God even one in me as I am in the Father as the Father is one in me that we may be one Behold Verily Verily I say unto my Servent Sidney I have looked upon thee \& thy works I have heard thy prayers \& prepared thee for a greater work thou art blessed for thou shalt do great things Behold thou wast sent forth even as John to prepare the way before me \& Elijah which should come \& thou knew it not thou didst Baptize by water unto Repentance but they received not the Holy Ghost but now I give unto you a commandment that thou shalt Baptize by water \& \sout{give} <\& they shall receive> the Holy Ghost by <the> laying on of hands even as the Apostles of old \& it shall come to pass that there shall be a great work in the land even among the gentiles for their folly their abominations shall be made manifest in the eyes of all People for I am God \& mine arm is <not> shortened \& I will shew miricles signs \& wonders unto all those who believe \sout{who believe} on my name \& whoso shall ask it in my name in faith they shall cast out Devils they shall heal the sick they shall cause the blind to receive their sight \& the deaf \sout{the} to hear \& the dumb to speak \& the lame to walk \& the time speedily cometh that great things are to be shewn forth unto the Children of men but without faith not any thing shall be shewn forth except desolations upon Babylon the same which has made all Nations drink of the wine of the wrath of her fornication \& there are none that do good except they that are ready to receive the fulness of my Gospel which I have sent forth to this generation Wherefore I have called upon the weak things \sout{of the world} they that are unlearned \& dispised to thresh the Nations by the Power of my spirit \& their arm shall be mine arm \& I will be their shield \& their Buckler \& I will gird up their loins \& they shall fight manfully for me \& their enemies shall be under their feet \& I will let fall the sword in their behalf \& by the fire of mine indignation will I preserve them \& the poor \& the meek shall have the Gospel preached unto them \& they shall be looking forth for the time of my coming for it is nigh at hand \& they shall learn the Parible of the figg tree for even now already summer is nigh \& I have sent forth the fullness of \sout{the} <my> Gospel by the hand of my servent Joseph \& in weakness have I blessed him \& I have given unto him the Keys of the mystery of those things which have been sealed even things which was from the foundation of the world \& the things which shall come from this time until the time of my coming if he abide in me \& if not another will I plant in his stead Wherefore watch over him that his faith fail not \& it shall be given by the comforter (the Holy Ghost) Which knoweth all things \& a commandment I give unto you that thou shalt write for him \& the scriptures shall be given even as they are in mine own bosom to the salvation of mine own elect for th[e]y will hear my voice \& shall see me \& shall not be asleep \& shall abide the day of my coming for they shall be purified even as I am pure \& now I Say unto you tarry with him \& he shall Journey with thee forsake him not \& shurely these things shall be fulfilled \& in as much as ye do not write behold it shall be given him to prophecy \& thou shalt Preach my gospel \& call on the Holy Prophets to prove his words as they shall be given him keep all the commandments \& covenants by which ye are bound \& I will cause the Heavens to shake for your Good \& satan shall tremble \& Zion shall rejoice upon the Hills \& florish \& Israel shall be saved in mine own due time \& by the Keys which have shall been given shall they be led \& no more be confounded at all[.] lift up your hearts \& be Glad your redemption draweth nigh fear not little flock the Kingdom is yours untill I come Behold I come quickly even so amen
				
		% -----
		\subsubsection{Revelation, 9 December 1830 [D\&C 36]}
		% -----
			\label{EMS-1830-JS-9-DEC-1830}
			% \label{EMS-1830-JS-DAY-3LETTERMONTH-YEAR}
		
			\begin{description}
				\item[Print Sources]
			\end{description}

				\begin{tabular}{p{0.5in}p{1in}p{0.5in}p{1in}}
				JSP	 	& D1:224--225	& JSR 		& 97	 \\
				DC	 	& Section 36	& HC 		& 1:131 \\
				HDDC 	& 467--483		& TCHC 		& 1:100 \\
				\end{tabular}
				% HDDC = woodford's DC diss; JSR = Marquardt's JS Revelations; TCHC = Vogel HC
			
			\begin{description}
				\item[Online Access] \href{https://www.josephsmithpapers.org/paper-summary/revelation-9-december-1830-dc-36/1}{Here}
				% \item[Analysis] \hyperref[DC36]{Here}
			\end{description}
			
			\begin{description}
				\item[Background]
			\end{description}
			
				% It might also be useful to give a two sentence context of the period: e.g., "This speech was given in Nauvoo to a small congregation one month prior to the destruction of the Nauvoo Expositor. These and other tensions may have been on the prophet's mind when he gave the speech."
			
			\begin{description}
				\item[Original Source]
			\end{description}
			
				% brief description of original source material, with reference to JSP or somewhere else for more info
				% Background material: it might be helpful to have a note that says whether a given document was presented as a speech, was written in a journal, etc.
				% Also a comment here about how reliable the source might be, or what shape it is in, if there are large omissions or parts that are hard to understand, and so on.
			
			\begin{description}
				\item[Text]
			\end{description}
			
				\begin{tightcenter}
				39\supersu{th}\supers{.} Commandment Dec\supersu{m}\supers{.} 9\supersu{th}\supers{.} AD 1830
				\end{tightcenter}
				
				A Commandment to Edard [Edward Partridge] his call to the Ministery \&c
				
				Saying thus saith the Lord God the mighty one of Israel behold I say unto you my Servent Edward thou art blessed \& thy sins are forgiven thee \& thou art called to preach my Gospel as with the voice of a Trump \& I will lay my hand upon you by the hand of my Servent sidney [Rigdon] \& thou shalt Receive my spirit (the Holy Ghost) even the comforter) which shall teach you the peacible things of the Kingdom \& thou shalt declare it with a loud voice Saying Blessed be the name of the most high God---
				
				And now this calling \& commandment give I unto all men that as many as shall come before my Servent Sidney \& Joseph embracing this calling \& commandment shall be ordained \& sent forth to preach the everlasting gospel among the Nation crying Repentance saying save yourselves from this untoward generation \& come forth out of the fire hating even the garment spotted with the flesh--- And this commandment shall be given unto the Elders of my Church that every man which will embrace it with singleness of heart may be ordained \& sent forth even as I have spoken I am Jesus Christ the Son of God Wherefore gird up your loins \& I will suddenly come to my temple even so amen
				
		% -----
		\subsubsection{Revelation, 30 December 1830 [D\&C 37]}
		% -----
			\label{EMS-1830-JS-30-DEC-1830}
			% \label{EMS-1830-JS-DAY-3LETTERMONTH-YEAR}
		
			\begin{description}
				\item[Print Sources]
			\end{description}

				\begin{tabular}{p{0.5in}p{1in}p{0.5in}p{1in}}
				JSP	 	& D1:226--227	& JSR 		& 97--98 \\
				DC	 	& Section 37	& HC 		& 1:139 \\
				HDDC 	& 484--489		& TCHC 		& 1:104 \\
				\end{tabular}
				% HDDC = woodford's DC diss; JSR = Marquardt's JS Revelations; TCHC = Vogel HC
			
			\begin{description}
				\item[Online Access] \href{https://www.josephsmithpapers.org/paper-summary/revelation-30-december-1830-dc-37/1}{Here}
				% \item[Analysis] \hyperref[DC37]{Here}
			\end{description}
			
			\begin{description}
				\item[Background]
			\end{description}
			
				% It might also be useful to give a two sentence context of the period: e.g., "This speech was given in Nauvoo to a small congregation one month prior to the destruction of the Nauvoo Expositor. These and other tensions may have been on the prophet's mind when he gave the speech."
			
			\begin{description}
				\item[Original Source]
			\end{description}
			
				% brief description of original source material, with reference to JSP or somewhere else for more info
				% Background material: it might be helpful to have a note that says whether a given document was presented as a speech, was written in a journal, etc.
				% Also a comment here about how reliable the source might be, or what shape it is in, if there are large omissions or parts that are hard to understand, and so on.
			
			\begin{description}
				\item[Text]
			\end{description}
			
				\begin{tightcenter}
				40\supers{th} Commandment AD 1830
				\end{tightcenter}
				
				A Revelation to Sidney [Rigdon] \& Joseph at a\sout{t} time that they went from Fayette to Canandaigua to translate \&c given at Canandaigua Ontario County State of New York
				
				A Commandment to sidney \& Joseph saying Behold I say unto ye that it is not Expedient in me that ye should Translate any more until ye shall go to the Ohio \& this because of the enemy \& for your sakes \& again I say unto you that ye shall not go untill ye have Preached my Gospel in those parts \& have strengthened up the Church whithersoever it is found \& more especially in Colesville for Behold they pray unto me in much faith \& again a commandment I give unto the Church that it is expedient in me that they should assemble together at the Ohio \sout{by} <against> the time that my Servent Oliver [Cowdery] shall return unto them Behold here is wisdom \& let evry man Choose for himself until I come amen even so amen
				
		% -----
		\subsubsection{Explanation of Scripture, 1830 [D\&C 74]}
		% -----
			\label{EMS-1830-JS-1830}
			% \label{EMS-1830-JS-DAY-3LETTERMONTH-YEAR}
		
			\begin{description}
				\item[Print Sources]
			\end{description}

				\begin{tabular}{p{0.5in}p{1in}p{0.5in}p{1in}}
				JSP	 	& D1:227--229	& JSR 		& 194--195 \\
				DC	 	& Section 74	& HC 		& 1:242 \\
				HDDC 	& 908--913		& TCHC 		& 1:168 \\
				\end{tabular}
				% HDDC = woodford's DC diss; JSR = Marquardt's JS Revelations; TCHC = Vogel HC
			
			\begin{description}
				\item[Online Access] \href{https://www.josephsmithpapers.org/paper-summary/explanation-of-scripture-1830-dc-74/1}{Here}
				% \item[Analysis] \hyperref[DC74]{Here}
			\end{description}
			
			\begin{description}
				\item[Background]
			\end{description}
			
				For the date of this document, see my notes at \hyperref[DC74]{section 74}.
			
				% It might also be useful to give a two sentence context of the period: e.g., "This speech was given in Nauvoo to a small congregation one month prior to the destruction of the Nauvoo Expositor. These and other tensions may have been on the prophet's mind when he gave the speech."
			
			\begin{description}
				\item[Original Source]
			\end{description}
			
				% brief description of original source material, with reference to JSP or somewhere else for more info
				% Background material: it might be helpful to have a note that says whether a given document was presented as a speech, was written in a journal, etc.
				% Also a comment here about how reliable the source might be, or what shape it is in, if there are large omissions or parts that are hard to understand, and so on.
			
			\begin{description}
				\item[Text]
			\end{description}
			
				An explanation of the \sout{first} Epistle to the first Corinthians 7 Chapter \& 14\supers{th} verse given to Joseph the Seer at Wayne County. N.Y 1830
				
				For the unbelieveing \sout{wife} husband is sanctified by the wife \& the unbelieveing wife is sanctified by the husband else were your Children unclean but now ere [are] they holy
				
				Now in the days of the Apostles the law of circumcision was had among \sout{them} all the Jews which believed not the Gospel of Jesus Christ, \& it came to pass that there arose a great contention among the People concerning the law of circumcision for the unbelieveing husband was desirous that his children should be circumcised \& become subject to the law of Moses which law was Fulfilled \& it came to pass that the Children being brought up in subjection to the law of moses \& gave heed to the traditions of their Fathers \& believed not the Gospel of Christ wherein they became unholy wherefore for this cause the Apostle wrote unto the Church giving unto them a commandment not of the Lord but of himself that a believer should not be united to an unbeliever except the law of Moses should be done away among them that their Children might remain without circumcision \& that the tradition might be done away which saith that little children are unholy for it was had among the Jews but little children are holy being sanctified through the atonement of Jesus Christ \& this is wat these scriptures mean
			
% -----------------------------------------------------------
\pagebreak{}
% -----------------------------------------------------------